%!TeX root=MemoriaTFG.tex

\chapter{Resultados}
\label{cap:resultados}

Este capítulo presenta los resultados experimentales obtenidos
durante la evaluación de modelos fundacionales dermatológicos y de
las contribuciones originales desarrolladas en este trabajo. Los
experimentos persiguen cuatro objetivos complementarios: evaluar la
capacidad de transferencia de los modelos sobre datasets externos
independientes; comparar distintas arquitecturas en tareas de
clasificación, segmentación y razonamiento multimodal; analizar
aspectos de robustez, eficiencia en etiquetas, equidad entre
fototipos y seguridad diagnóstica; y, finalmente, validar la
transferencia de estas representaciones a un entorno clínico en
castellano mediante DermapixelAI~1.0 y explorar su adaptación
eficiente a través de Dermapixel R0. Estos cuatro objetivos
concretan los objetivos formales del trabajo (sección~\ref{sec:objetivos}).

La secuencia experimental se organiza de forma acumulativa. En
primer lugar, se verifica la independencia entre los conjuntos de
evaluación y los datos utilizados durante el preentrenamiento, lo
que garantiza la validez de las comparaciones posteriores
(sección~\ref{sec:auditoria-resultados}). Sobre esta base se estudia
la transferencia supervisada mediante sondeo lineal, ajuste fino y
eficiencia en etiquetas; posteriormente se analizan tareas
específicas de segmentación, clasificación \emph{zero-shot} y
razonamiento multimodal; a continuación se examinan aspectos de
equidad y seguridad diagnóstica; y, finalmente, se presentan las
contribuciones propias del trabajo sobre DermapixelAI~1.0 y
Dermapixel R0, junto con la interpretabilidad mediante \emph{Sparse
Autoencoders} y una síntesis comparativa.

En términos generales, los resultados muestran que los modelos
especializados en dermatología superan de forma consistente a las
alternativas generalistas, que la adaptación eficiente permite
transferir estas capacidades a un entorno clínico en castellano, y
que la calidad de los datos y de las representaciones resulta más
determinante que el simple incremento del tamaño de los modelos. La
figura~\ref{fig:cap4-estructura} resume esta organización.

\begin{figure}[H]
\centering
\small
\renewcommand{\arraystretch}{1.1}
\begin{tabular}{c}
\fbox{\textbf{Validez:} auditoría de solapamiento con Derm1M (§\ref{sec:auditoria-resultados})} \\
$\downarrow$ \\
\fbox{\textbf{Transferencia:} sondeo lineal, ajuste fino y eficiencia en etiquetas (datasets externos)} \\
$\downarrow$ \\
\fbox{\textbf{Segmentación} promptable basada en SAM/SAM2 (§\ref{sec:seg-results})} \\
$\downarrow$ \\
\fbox{\textbf{Zero-shot y modelos de lenguaje multimodales}} \\
$\downarrow$ \\
\fbox{\textbf{Equidad por fototipo y seguridad diagnóstica}} \\
$\downarrow$ \\
\fbox{\textbf{DermapixelAI~1.0 y Dermapixel R0} (validación en castellano y adaptación, O2/O3)} \\
$\downarrow$ \\
\fbox{\textbf{Interpretabilidad (SAE) y síntesis comparativa}} \\
\end{tabular}
\caption{Organización experimental del capítulo: de la validación de
la independencia de los datos a las contribuciones propias en
castellano.}
\label{fig:cap4-estructura}
\end{figure}

\section{Validez experimental: auditoría de solapamiento con Derm1M}
\label{sec:auditoria-resultados}

La validez de toda comparación de transferencia y \emph{zero-shot}
de este capítulo depende de que los conjuntos de test no se solapen
con el dataset de preentrenamiento; por ello el capítulo abre
comprobándolo. La auditoría intersecta los hashes MD5 de los doce
datasets externos con los $415\,451$ hashes únicos de la fracción
accesible de Derm1M e identifica un único dataset con solapamiento
exacto: \textbf{Dermnet}, cuyas $18\,302$ imágenes están todas
contenidas en Derm1M ($100\%$). Los once datasets restantes no
arrojan ninguna coincidencia. En particular, HIBA y MSKCC ---cuyo
origen en el ISIC archive motivaba una sospecha de procedencia---
presentan $0$ coincidencias MD5, por lo que esa sospecha no se
traduce en solapamiento exacto. El resultado completo por dataset
figura en el Anexo~\ref{cap:anexoa}.

El solapamiento de Dermnet no invalida las cifras que se reportan en
la sección~\ref{sec:rep-panderm} para PanDerm, DINOv2,
ConvNeXt-Large, EfficientNetV2-Large y SigLIP-Large, cuyo
preentrenamiento no incluye Derm1M. Sí condiciona, en cambio, la
interpretación de las cifras zero-shot de la familia DermLIP sobre
Dermnet (sección~\ref{sec:zs-results}), ya que DermLIP se entrenó
sobre Derm1M y vio durante el preentrenamiento la totalidad de ese
conjunto.

\section{Reproducción de PanDerm: sondeo lineal sobre diez datasets}
\label{sec:rep-panderm}

¿Bastan las representaciones de un modelo fundacional dermatológico,
sin reajustar el encoder, para clasificar lesiones en datasets que
nunca vio? Esta pregunta determina si el preentrenamiento de dominio
aporta valor reutilizable o exige reentrenar en cada despliegue. La
evidencia muestra que el codificador de mayor capacidad transfiere
de forma robusta entre datasets y que su ventaja aumenta
especialmente en los escenarios con mayor heterogeneidad de
adquisición y complejidad taxonómica; los resultados sobre diez
datasets externos lo cuantifican a continuación.

La tabla~\ref{tab:lp-panderm} aplica este protocolo a PanDerm Base y
Large sobre los diez datasets externos contrastados, conforme a
sección~\ref{sec:lp}. Se reportan cinco métricas: \emph{accuracy}
(Acc), \emph{balanced accuracy} (BAcc), AUROC \emph{one-vs-rest}, F1
ponderada (W-F1) y coeficiente kappa de Cohen.

\begin{table}[H]
\centering
\footnotesize
\setlength{\tabcolsep}{4pt}
\renewcommand{\arraystretch}{0.95}
\caption{Sondeo lineal con PanDerm Base y Large sobre diez datasets
externos. El asterisco (\textsuperscript{$\ast$}) marca Dermnet, el
único dataset con solapamiento exacto frente a Derm1M ($100\%$ por
hash MD5; sección~\ref{sec:auditoria-resultados}).
Mejores valores por dataset y métrica en negrita.}
\label{tab:lp-panderm}
\begin{tabular}{llrrrrr}
\hline
Dataset       & Modelo & Acc            & BAcc           & AUROC          & W-F1           & Kappa \\
\hline
HAM10000      & Base   & $0{,}853$      & $0{,}381$      & $0{,}892$      & $0{,}829$      & $0{,}709$ \\
HAM10000      & Large  & $\mathbf{0{,}888}$ & $\mathbf{0{,}575}$ & $\mathbf{0{,}954}$ & $\mathbf{0{,}880}$ & $\mathbf{0{,}814}$ \\
BCN20000      & Base   & $0{,}659$      & $0{,}292$      & $0{,}855$      & $0{,}615$      & $0{,}447$ \\
BCN20000      & Large  & $\mathbf{0{,}702}$ & $\mathbf{0{,}382}$ & $\mathbf{0{,}903}$ & $\mathbf{0{,}676}$ & $\mathbf{0{,}527}$ \\
PAD-UFES-20   & Base   & $0{,}725$      & $0{,}510$      & $0{,}913$      & $0{,}692$      & $0{,}369$ \\
PAD-UFES-20   & Large  & $\mathbf{0{,}772}$ & $\mathbf{0{,}642}$ & $\mathbf{0{,}949}$ & $\mathbf{0{,}760}$ & $\mathbf{0{,}527}$ \\
Dermnet$^\ast$& Base   & $0{,}450$      & $0{,}381$      & $0{,}888$      & $0{,}432$      & $0{,}437$ \\
Dermnet$^\ast$& Large  & $\mathbf{0{,}550}$ & $\mathbf{0{,}495}$ & $\mathbf{0{,}931}$ & $\mathbf{0{,}540}$ & $\mathbf{0{,}555}$ \\
WSI patches   & Base   & $0{,}781$      & $0{,}640$      & $0{,}976$      & $0{,}774$      & $0{,}842$ \\
WSI patches   & Large  & $\mathbf{0{,}868}$ & $\mathbf{0{,}792}$ & $\mathbf{0{,}991}$ & $\mathbf{0{,}871}$ & $\mathbf{0{,}896}$ \\
DDI           & Base   & $\mathbf{0{,}847}$ & $0{,}714$      & $0{,}827$      & $0{,}836$      & $0{,}482$ \\
DDI           & Large  & $\mathbf{0{,}847}$ & $\mathbf{0{,}764}$ & $\mathbf{0{,}860}$ & $\mathbf{0{,}846}$ & $\mathbf{0{,}535}$ \\
Derm7pt clín. & Base   & $0{,}762$      & $0{,}675$      & $0{,}808$      & $0{,}736$      & $0{,}397$ \\
Derm7pt clín. & Large  & $\mathbf{0{,}798}$ & $\mathbf{0{,}732}$ & $\mathbf{0{,}869}$ & $\mathbf{0{,}784}$ & $\mathbf{0{,}506}$ \\
Derm7pt dermo & Base   & $0{,}818$      & $0{,}749$      & $0{,}841$      & $0{,}811$      & $0{,}528$ \\
Derm7pt dermo & Large  & $\mathbf{0{,}836}$ & $\mathbf{0{,}766}$ & $\mathbf{0{,}858}$ & $\mathbf{0{,}829}$ & $\mathbf{0{,}570}$ \\
HIBA          & Base   & $0{,}883$      & $0{,}595$      & $0{,}894$      & $0{,}853$      & $0{,}275$ \\
HIBA          & Large  & $\mathbf{0{,}904}$ & $\mathbf{0{,}701}$ & $\mathbf{0{,}942}$ & $\mathbf{0{,}892}$ & $\mathbf{0{,}494}$ \\
MSKCC         & Base   & $0{,}721$      & $\mathbf{0{,}654}$ & $0{,}746$      & $0{,}720$      & $0{,}309$ \\
MSKCC         & Large  & $\mathbf{0{,}746}$ & $0{,}644$      & $\mathbf{0{,}751}$ & $\mathbf{0{,}732}$ & $\mathbf{0{,}316}$ \\
\hline
\emph{Media}  & Base   & $0{,}750$      & $0{,}559$      & $0{,}865$      & $0{,}730$      & --- \\
\emph{Media}  & Large  & $\mathbf{0{,}791}$ & $\mathbf{0{,}649}$ & $\mathbf{0{,}901}$ & $\mathbf{0{,}781}$ & --- \\
\hline
\end{tabular}
\end{table}

PanDerm Large supera a PanDerm Base en accuracy, AUROC y F1
ponderada sobre los diez datasets evaluados. La mejora media
absoluta alcanza $+4{,}1$\,pp en accuracy, $+9{,}0$\,pp en
\emph{balanced accuracy} y $+3{,}6$\,pp en AUROC
(figura~\ref{fig:lp-per-dataset}). La única excepción aparece en la
BAcc de MSKCC, donde PanDerm Base obtiene una ventaja marginal de
$1{,}0$\,pp ($0{,}654$ frente a $0{,}644$).

La magnitud de la mejora no es uniforme entre datasets. Las
diferencias más amplias se observan en Dermnet ($+10{,}0$\,pp
accuracy), WSI~patches ($+8{,}7$\,pp) y PAD-UFES-20 ($+4{,}7$\,pp),
conjuntos caracterizados por una mayor diversidad de clases,
modalidades o condiciones de adquisición. Este patrón sugiere que el
incremento de capacidad del encoder resulta especialmente
beneficioso cuando aumenta la complejidad del problema de
generalización. La modalidad con mejor rendimiento absoluto es la
histopatología (WSI~patches), donde PanDerm Large alcanza una
accuracy de $0{,}868$ y una AUROC de $0{,}991$. Dermnet queda
marcado con asterisco por estar íntegramente contenido en Derm1M
($100\%$ por hash MD5); su interpretación se discute en
sección~\ref{sec:auditoria-resultados}.

\begin{figure}[H]
\centering
\includegraphics[width=0.92\textwidth]{fig_lp_per_dataset.png}
\caption{Sondeo lineal con PanDerm Base y Large sobre los diez
datasets externos contrastados. Diferencia absoluta en BAcc
($+9{,}0$\,pp de promedio) y AUROC ($+3{,}6$\,pp de
promedio) a favor de PanDerm Large.}
\label{fig:lp-per-dataset}
\end{figure}

En conjunto, estos resultados indican que el beneficio del
preentrenamiento dermatológico no se limita al dataset de origen,
sino que se transfiere de forma consistente a múltiples cohortes
externas, modalidades de imagen y condiciones de adquisición.

\section{Benchmark comparativo entre familias}
\label{sec:benchmark}

¿La ventaja observada para PanDerm Large refleja realmente un
beneficio del preentrenamiento dermatológico o simplemente una
comparación favorable frente a su versión Base? Para responder a
esta cuestión resulta necesario situarlo frente a familias de
modelos con filosofías de aprendizaje diferentes: arquitecturas
convolucionales supervisadas, modelos auto-supervisados generalistas
y sistemas visión-lenguaje entrenados a gran escala. La
tabla~\ref{tab:benchmark} compara cinco modelos representativos de
estas familias sobre HAM10000, PAD-UFES-20 y DDI bajo el mismo
protocolo de sondeo lineal.

\begin{table}[H]
\centering
\footnotesize
\setlength{\tabcolsep}{4pt}
\renewcommand{\arraystretch}{0.95}
\caption{Benchmark comparativo bajo sondeo lineal sobre HAM10000,
PAD-UFES-20 y DDI. Mejor valor por columna en negrita.}
\label{tab:benchmark}
\adjustbox{max width=\textwidth}{%
\begin{tabular}{lrrrrrrr}
\hline
Modelo                  & Tipo          & Params & HAM Acc       & HAM AUROC     & PAD Acc       & PAD AUROC     & DDI AUROC \\
\hline
ConvNeXt-Large          & ConvNet sup.  & $197$\,M & $0{,}883$ & $\mathbf{0{,}967}$ & $0{,}668$ & $0{,}894$ & $0{,}774$ \\
EfficientNetV2-Large    & ConvNet sup.  & $118$\,M & $0{,}860$ & $0{,}950$ & $0{,}690$ & $0{,}895$ & $0{,}766$ \\
DINOv2 ViT-L/14         & SSL general   & $307$\,M & $0{,}744$ & $0{,}847$ & ---       & ---       & --- \\
PanDerm Large           & SSL derm.     & $307$\,M & $0{,}888$ & $0{,}954$ & $\mathbf{0{,}772}$ & $\mathbf{0{,}949}$ & $\mathbf{0{,}860}$ \\
SigLIP-Large SO400M     & Vis-leng.\ gen.& $400$\,M & $\mathbf{0{,}900}$ & $0{,}971$ & $0{,}718$ & $0{,}903$ & $0{,}793$ \\
\hline
\end{tabular}}
\end{table}

La tabla~\ref{tab:gap-cnn} cuantifica la ventaja de PanDerm Large
sobre la mejor CNN supervisada en cada dataset.

\begin{table}[H]
\centering
\footnotesize
\setlength{\tabcolsep}{4pt}
\caption{PanDerm Large frente a la mejor CNN supervisada (la mejor de
ConvNeXt-L y EfficientNetV2-L para cada métrica).}
\label{tab:gap-cnn}
\begin{tabular}{lrrrrrr}
\hline
Dataset     & CNN Acc       & PanDerm Acc   & $\Delta$Acc (pp) & CNN AUROC     & PanDerm AUROC & $\Delta$AUROC (pp) \\
\hline
HAM10000    & $0{,}883$     & $0{,}888$     & $+0{,}5$         & $0{,}967$     & $0{,}954$     & $-1{,}3$ \\
PAD-UFES-20 & $0{,}690$     & $0{,}772$     & $+8{,}2$         & $0{,}895$     & $0{,}949$     & $+5{,}4$ \\
DDI         & $0{,}818$     & $0{,}847$     & $+2{,}9$         & $0{,}774$     & $0{,}860$     & $+8{,}6$ \\
\hline
\end{tabular}
\end{table}

PanDerm Large supera a la mejor CNN supervisada en accuracy sobre
los tres datasets (figura~\ref{fig:lp-benchmark}). La magnitud de la
ventaja no es uniforme entre datasets: $+0{,}5$\,pp en HAM10000
(dermatoscopia estandarizada), frente a $+8{,}2$\,pp en PAD-UFES-20
(fotografía clínica con dispositivo móvil) y $+2{,}9$\,pp en DDI
(fotografía clínica con representación equilibrada de fototipos). El
patrón en AUROC es similar, salvo en HAM10000, donde ConvNeXt-Large
($0{,}967$) supera a PanDerm Large por $1{,}3$\,pp. SigLIP-Large
SO400M, vision-lenguaje generalista, logra la mejor accuracy y AUROC
en HAM10000 ($0{,}900 / 0{,}971$). Una posible explicación es la
proximidad entre las imágenes dermatoscópicas estandarizadas de este
dataset y parte del contenido visual presente en su preentrenamiento
a gran escala, aunque esta hipótesis no puede verificarse
directamente a partir de los datos disponibles. PanDerm Large, en
cambio, domina en PAD-UFES-20 y DDI, los datasets clínicos más
alejados de los datasets web generalistas.

\begin{figure}[H]
\centering
\includegraphics[width=0.88\textwidth]{fig_lp_comparison.png}
\caption{Benchmark comparativo bajo sondeo lineal sobre HAM10000,
PAD-UFES-20 y DDI. ConvNeXt-Large supera a PanDerm Large en AUROC
sobre HAM10000 ($-1{,}3$\,pp); PanDerm Large domina sobre
PAD-UFES-20 ($+8{,}2$\,pp Acc) y DDI ($+8{,}6$\,pp AUROC).}
\label{fig:lp-benchmark}
\end{figure}

En conjunto, los resultados muestran que el preentrenamiento
específico en dermatología ofrece una ventaja especialmente evidente
cuando aumenta la heterogeneidad del dominio de evaluación. Aunque
modelos generalistas de gran escala como SigLIP pueden igualar o
incluso superar a PanDerm Large en escenarios dermatoscópicos muy
estandarizados, los modelos especializados mantienen una capacidad
de generalización más robusta sobre datasets clínicos externos y
distribuciones alejadas del dominio web.

\section{Ajuste fino supervisado}
\label{sec:ft-results}

Si la representación congelada ya transfiere, ¿qué se gana al
descongelar y reajustar el encoder? La respuesta delimita cuándo
conviene el ajuste fino y cuándo es contraproducente, una decisión
central para los datasets clínicos pequeños. En los datasets
evaluados, el ajuste fino mejora claramente cuando existe una
cantidad moderada o elevada de datos de entrenamiento y puede
degradar el rendimiento en escenarios de baja disponibilidad de
muestras. La tabla~\ref{tab:ft-panderm} lo concreta comparando el
sondeo lineal con el ajuste fino supervisado de PanDerm Base sobre
cuatro datasets de tamaño dispar, con la configuración descrita
en sección~\ref{sec:ft}. Sobre HAM10000 se añade además
\emph{test-time augmentation} con cinco augmentaciones
deterministas.

\begin{table}[H]
\centering
\footnotesize
\setlength{\tabcolsep}{4pt}
\renewcommand{\arraystretch}{0.95}
\caption{Sondeo lineal (LP) frente a ajuste fino (FT) con PanDerm
Base. $\Delta$BAcc en puntos porcentuales respecto al LP.}
\label{tab:ft-panderm}
\begin{tabular}{llrrrrr}
\hline
Dataset      & Modo & Acc           & BAcc          & AUROC         & W-F1          & $\Delta$BAcc \\
\hline
HAM10000     & LP   & $0{,}853$     & $0{,}381$     & $0{,}892$     & $0{,}829$     & --- \\
HAM10000     & FT   & $\mathbf{0{,}920}$ & $\mathbf{0{,}852}$ & $\mathbf{0{,}978}$ & $\mathbf{0{,}923}$ & $+47{,}1$ \\
PAD-UFES-20  & LP   & $0{,}725$     & $0{,}510$     & $0{,}913$     & $0{,}692$     & --- \\
PAD-UFES-20  & FT   & $\mathbf{0{,}755}$ & $\mathbf{0{,}704}$ & $\mathbf{0{,}935}$ & $\mathbf{0{,}757}$ & $+19{,}4$ \\
DDI          & LP   & $\mathbf{0{,}847}$ & $\mathbf{0{,}714}$ & $\mathbf{0{,}827}$ & $\mathbf{0{,}836}$ & --- \\
DDI          & FT   & $0{,}774$     & $0{,}693$     & $0{,}682$     & $0{,}780$     & $-2{,}1$ \\
MSKCC        & LP   & $0{,}721$     & $0{,}654$     & $\mathbf{0{,}746}$ & $0{,}720$     & --- \\
MSKCC        & FT   & $\mathbf{0{,}731}$ & $\mathbf{0{,}684}$ & $0{,}719$     & $\mathbf{0{,}735}$ & $+3{,}0$ \\
\hline
\end{tabular}
\end{table}

El ajuste fino eleva la BAcc en HAM10000 ($+47{,}1$\,pp) y
PAD-UFES-20 ($+19{,}4$\,pp). La mejora observada se concentra
principalmente en las clases minoritarias (melanoma, \emph{akiec},
\emph{df}, \emph{vasc} en HAM10000), responsables de la baja
\emph{balanced accuracy} del sondeo lineal: el ajuste fino adapta
las representaciones visuales a estas categorías poco frecuentes y
reduce parcialmente el desequilibrio. El efecto se invierte cuando
escasean los datos. Sobre DDI, con $N_{\mathrm{train}} = 427$, el
sondeo lineal supera al ajuste fino en todas las métricas, señal de
sobreajuste del encoder pese a las augmentaciones. Sobre MSKCC
($N_{\mathrm{train}} = 6\,189$), la mejora de BAcc es marginal
($+3{,}0$\,pp) y la AUROC baja un poco. Estos resultados sugieren
que, en escenarios con pocos cientos de imágenes de entrenamiento,
el sondeo lineal sobre las representaciones congeladas puede
resultar más robusto que el ajuste fino del encoder
(figura~\ref{fig:ft-comparison}).

\begin{figure}[H]
\centering
\includegraphics[width=0.85\textwidth]{fig_ft_comparison.png}
\caption{Comparación entre sondeo lineal y ajuste fino con
PanDerm Base sobre cuatro datasets. El ajuste fino aporta
mejoras sustanciales sobre HAM10000 y PAD-UFES-20
($+47{,}1$\,pp y $+19{,}4$\,pp BAcc respectivamente) pero
degrada el rendimiento sobre DDI ($N_{\mathrm{train}} = 427$),
caso compatible con sobreajuste pese a las augmentaciones.}
\label{fig:ft-comparison}
\end{figure}

En conjunto, los resultados indican que la utilidad del ajuste fino
depende de la disponibilidad de datos etiquetados. Cuando el
conjunto de entrenamiento es suficientemente grande, la adaptación
supervisada del encoder produce mejoras sustanciales respecto al
sondeo lineal; sin embargo, esta ventaja desaparece e incluso puede
invertirse en escenarios con disponibilidad limitada de muestras,
donde las representaciones fundacionales congeladas ofrecen una
alternativa más estable. La siguiente cuestión es cuántas etiquetas
son realmente necesarias para aprovechar estas representaciones
fundacionales.

\section{Eficiencia en etiquetas}
\label{sec:label-eff-results}

Para responder a esta pregunta, se mide el rendimiento de PanDerm
Large bajo submuestreos estratificados crecientes del conjunto de
entrenamiento. El etiquetado experto es el principal cuello de
botella en dermatología, por lo que la respuesta tiene impacto
práctico directo. El hallazgo es que una fracción mínima del
conjunto basta para casi saturar el rendimiento. La
tabla~\ref{tab:label-eff} lo demuestra con submuestreos al $1\%$,
$5\%$, $10\%$, $20\%$, $50\%$ y $100\%$ sobre HAM10000 y
PAD-UFES-20.

\begin{table}[H]
\centering
\footnotesize
\setlength{\tabcolsep}{4pt}
\renewcommand{\arraystretch}{0.95}
\caption{Eficiencia en etiquetas (sondeo lineal con PanDerm Large).
$N_{\mathrm{train}}$: tamaño efectivo del subconjunto de
entrenamiento. Mejor valor por columna en negrita.}
\label{tab:label-eff}
\begin{tabular}{lrrrrr}
\hline
$\%$ entrenam.  & $N_{\mathrm{train}}$ & Acc           & BAcc          & AUROC         & W-F1 \\
\hline
\multicolumn{6}{l}{\textit{HAM10000}} \\
$1\%$           & $82$                 & $0{,}850$     & $0{,}406$     & $0{,}918$     & $0{,}828$ \\
$5\%$           & $410$                & $0{,}870$     & $0{,}481$     & $0{,}932$     & $0{,}855$ \\
$10\%$          & $820$                & $0{,}875$     & $0{,}480$     & $0{,}939$     & $0{,}863$ \\
$20\%$          & $1\,641$             & $0{,}882$     & $0{,}547$     & $0{,}951$     & $0{,}873$ \\
$50\%$          & $4\,103$             & $\mathbf{0{,}894}$ & $\mathbf{0{,}589}$ & $0{,}952$ & $\mathbf{0{,}887}$ \\
$100\%$         & $8\,207$             & $0{,}888$     & $0{,}575$     & $\mathbf{0{,}954}$ & $0{,}880$ \\
\hline
\multicolumn{6}{l}{\textit{PAD-UFES-20}} \\
$1\%$           & $14$                 & $0{,}740$     & $0{,}599$     & $0{,}907$     & $0{,}727$ \\
$5\%$           & $74$                 & $0{,}751$     & $0{,}632$     & $0{,}919$     & $0{,}742$ \\
$10\%$          & $149$                & $0{,}738$     & $0{,}593$     & $0{,}923$     & $0{,}724$ \\
$20\%$          & $298$                & $0{,}748$     & $0{,}587$     & $0{,}931$     & $0{,}728$ \\
$50\%$          & $746$                & $0{,}757$     & $0{,}611$     & $0{,}941$     & $0{,}742$ \\
$100\%$         & $1\,493$             & $\mathbf{0{,}772}$ & $\mathbf{0{,}642}$ & $\mathbf{0{,}949}$ & $\mathbf{0{,}760}$ \\
\hline
\end{tabular}
\end{table}

La AUROC satura muy pronto (figura~\ref{fig:label-efficiency}). Sobre HAM10000, con solo el $1\%$ del
entrenamiento ($82$ imágenes), PanDerm Large recupera $0{,}850$ de
accuracy y $0{,}918$ de AUROC, a $3{,}8$\,pp y $3{,}6$\,pp del
rendimiento con el conjunto completo. Con el $5\%$ ($410$ imágenes)
la AUROC llega a $0{,}932$, el $97{,}7\%$ del máximo. A diferencia
de la AUROC, la \emph{balanced accuracy} continúa mostrando
sensibilidad al incremento del volumen de entrenamiento y alcanza su
mejor valor observado con el $50\%$ de los datos ($0{,}589$). El
patrón se repite sobre PAD-UFES-20: con $14$ imágenes de
entrenamiento el modelo obtiene $0{,}740$ de accuracy y $0{,}907$ de
AUROC, el $95{,}5\%$ de la AUROC con el conjunto completo. Este
patrón es coherente con la hipótesis de que gran parte del
conocimiento visual relevante ya se encuentra codificado en las
representaciones aprendidas durante el preentrenamiento; en
consecuencia, la cabeza lineal requiere relativamente pocos ejemplos
para alcanzar una fracción elevada del rendimiento máximo.

\begin{figure}[H]
\centering
\includegraphics[width=0.88\textwidth]{fig_label_efficiency.png}
\caption{Eficiencia en etiquetas con PanDerm Large sobre HAM10000
y PAD-UFES-20. La curva AUROC satura tempranamente: con
$1\%$ del conjunto de entrenamiento se recupera más del
$95\%$ del rendimiento alcanzable con el conjunto completo
en ambos datasets.}
\label{fig:label-efficiency}
\end{figure}

En conjunto, los resultados indican que las representaciones
aprendidas por PanDerm Large son altamente eficientes desde el punto
de vista supervisado. La mayor parte de la capacidad discriminativa
se encuentra ya presente en el encoder preentrenado, lo que reduce
sustancialmente la dependencia de grandes volúmenes de datos
etiquetados y resulta especialmente relevante para escenarios
clínicos donde el etiquetado experto constituye el principal cuello
de botella.

\section{Segmentación binaria de lesión}
\label{sec:seg-results}

Hasta aquí los protocolos atacaban clasificación. La segmentación
plantea otra pregunta: delimitar el contorno de la lesión píxel a
píxel.

La pregunta científica del bloque es si un modelo fundacional de
segmentación adaptado al dominio dermatológico alcanza una
delimitación de lesión de calidad clínica y qué factores la
determinan. La evaluación se articula en tres niveles. En primer
lugar se establece una referencia mediante la adaptación de
SAM2.1-Large al dominio dermatológico. A continuación se comparan
cinco arquitecturas \emph{promptables} bajo un protocolo homogéneo
para aislar el efecto de la arquitectura, el preentrenamiento médico
y el tamaño del \emph{backbone}. Finalmente, diversos análisis
complementarios estudian la dependencia del \emph{prompt} ---mediante
un control automático sin caja basado en una U-Net---, la
estabilidad del entrenamiento y el límite impuesto por la
variabilidad entre anotadores humanos.

La tabla~\ref{tab:seg} reporta la segmentación binaria sobre
ISIC2018 con SAM2.1-Large adaptado mediante \emph{fine-tuning} de su
decodificador, junto a tres referencias: SAM2 sin adaptación
(\emph{zero-shot}), \emph{Convolutional AutoEncoder} con backbone
PanDerm Large entrenado desde cero (CAE-seg) y la cifra publicada
por Yan et~al.~\cite{panderm2025}.

\begin{table}[H]
\centering
\footnotesize
\setlength{\tabcolsep}{4pt}
\renewcommand{\arraystretch}{0.95}
\caption{Segmentación binaria de lesión: Dice de test sobre tres
datasets dermatoscópicos. Modelos entrenados sobre ISIC2018 y
evaluados sin reentrenamiento sobre ISIC2017 y PH2.}
\label{tab:seg}
\begin{tabular}{lrrr}
\hline
Configuración                         & ISIC2018 Dice & ISIC2017 Dice & PH2 Dice \\
\hline
SAM2 \emph{zero-shot} (bbox de GT)    & ---           & $0{,}797$     & $0{,}886$ \\
CAE-seg (PanDerm-L) FT ISIC2018       & $0{,}894$     & $0{,}879$     & $0{,}924$ \\
PanDerm \cite{panderm2025} (ref.)     & $0{,}921$     & ---           & --- \\
\textbf{SAM2.1-L (decoder FT) ISIC2018}  & $\mathbf{0{,}947}$ & $\mathbf{0{,}945}$ & $\mathbf{0{,}960}$ \\
\hline
\end{tabular}
\end{table}

La tabla~\ref{tab:seg-detalle} amplía la información con la
desviación estándar y el IoU sobre los datasets de transferencia.

\begin{table}[H]
\centering
\footnotesize
\setlength{\tabcolsep}{4pt}
\caption{Generalización fuera de dominio: Dice e IoU (media $\pm$
desv.~estándar) sobre ISIC2017 ($N = 600$) y PH2 ($N = 200$).}
\label{tab:seg-detalle}
\begin{tabular}{llrr}
\hline
Modelo                              & Dataset      & Dice              & IoU \\
\hline
SAM2 \emph{zero-shot} (bbox)        & ISIC2017     & $0{,}797 \pm 0{,}158$ & $0{,}686 \pm 0{,}183$ \\
SAM2 \emph{zero-shot} (bbox)        & PH2          & $0{,}886 \pm 0{,}199$ & $0{,}830 \pm 0{,}198$ \\
SAM2.1-L (decoder FT) ISIC2018      & ISIC2017     & $0{,}945 \pm 0{,}046$ & $0{,}898 \pm 0{,}065$ \\
SAM2.1-L (decoder FT) ISIC2018      & PH2          & $0{,}960 \pm 0{,}022$ & $0{,}925 \pm 0{,}039$ \\
CAE-seg (PanDerm-L) FT ISIC2018     & ISIC2017     & $0{,}879 \pm 0{,}139$ & $0{,}803 \pm 0{,}166$ \\
CAE-seg (PanDerm-L) FT ISIC2018     & PH2          & $0{,}924 \pm 0{,}066$ & $0{,}865 \pm 0{,}098$ \\
\hline
\end{tabular}
\end{table}

SAM2.1-L con \emph{fine-tuning} del decodificador alcanza Dice
$0{,}947$ sobre ISIC2018. Supera en $+2{,}6$\,pp la cifra
publicada por Yan et~al.~\cite{panderm2025} ($0{,}921$) y en
$+5{,}3$\,pp al CAE-seg con backbone PanDerm-L. La transferencia
fuera de dominio es mínima: el Dice desciende solo $0{,}002$ sobre
ISIC2017 ($0{,}947 \to 0{,}945$) y, sobre PH2, sube hasta
$0{,}960$, lo que se explica por la mayor limpieza visual de PH2.
Adaptar el decodificador aporta frente a SAM2 sin adaptación
$+14{,}8$\,pp de Dice sobre ISIC2017 ($0{,}797 \to 0{,}945$) y
$+7{,}4$\,pp sobre PH2 ($0{,}886 \to 0{,}960$). Estas diferencias
cuantifican el aporte específico de la adaptación de dominio
(figura~\ref{fig:seg-generalization}).

\begin{figure}[H]
\centering
\includegraphics[width=0.88\textwidth]{fig_seg_generalization.png}
\caption{Segmentación binaria de lesión. SAM2.1-Large con
\emph{fine-tuning} del decodificador alcanza Dice $0{,}947$ sobre ISIC2018,
$+2{,}6$\,pp sobre el resultado publicado por
Yan~et~al.~\cite{panderm2025}, con transferencia OOD
prácticamente nula sobre ISIC2017 ($0{,}945$) y PH2
($0{,}960$).}
\label{fig:seg-generalization}
\end{figure}

\subsection{Comparativa de arquitecturas promptables}
\label{sec:seg-arquitecturas}

La cifra anterior mezcla tres factores: la generación de
arquitectura, el preentrenamiento de dominio y el tamaño de
\emph{backbone}. Para separarlos, se evalúan de forma pareada cinco
arquitecturas de segmentación promptable sobre un dataset
dermatoscópico ampliado y deduplicado por hash MD5 de píxel
($14\,201$ imágenes únicas de ISIC2018,
ISIC2017, HAM10000-seg y PH2). El régimen es común a todos los
modelos: \emph{fine-tuning} denso del decodificador y el
\emph{promptable encoder}, encoder visual congelado,
\emph{prompt} de \emph{bounding box} derivada de la máscara de
referencia, AdamW con $\eta = 1\times 10^{-4}$, pérdida
Dice$+$BCE y $30$ épocas con \emph{early stopping}. La evaluación
emplea el mismo test canónico de ISIC2018 ($N = 1\,000$,
blindado frente a entrenamiento) que la cifra de referencia
$0{,}947$ de la tabla~\ref{tab:seg}, lo que mantiene la
comparación bit a bit. La tabla~\ref{tab:seg-arquitecturas}
reporta el Dice resultante.

\begin{table}[H]
\centering
\footnotesize
\setlength{\tabcolsep}{4pt}
\renewcommand{\arraystretch}{0.95}
\caption{Comparativa de cinco arquitecturas promptables sobre el
test canónico de ISIC2018 ($N = 1\,000$). Régimen común de
\emph{fine-tuning} del decodificador con \emph{prompt} de
\emph{bounding box}. Las diferencias de Dice son modestas en valor
absoluto pero significativas al test de Wilcoxon
pareado~\cite{wilcoxon1945}
($p < 10^{-6}$ en todos los contrastes).}
\label{tab:seg-arquitecturas}
\begin{tabular}{llllrr}
\hline
Modelo            & Arquitectura       & Preentr. & Params & Dice & Latencia \\
\hline
MedSAM2-tiny~\cite{medsam2} & SAM2~\cite{sam2}   & médico   & $39$\,M  & $\mathbf{0{,}9556}$ & $26$\,ms \\
SAM2.1-tiny       & SAM2~\cite{sam2}   & general  & $39$\,M  & $0{,}9517$ & $26$\,ms \\
SAM2.1-Large      & SAM2~\cite{sam2}   & general  & $224$\,M & $0{,}9503$ & $96$\,ms \\
SAM-Med2D~\cite{sammed2d} & SAM~\cite{sam2023} & médico   & $271$\,M & $0{,}9465$ & $20$\,ms \\
MedSAM v1~\cite{medsam}   & SAM~\cite{sam2023} & médico   & $94$\,M  & $0{,}9458$ & $74$\,ms \\
\hline
\end{tabular}
\end{table}

La comparativa deja tres regularidades. Primero, la generación de
arquitectura es el factor dominante: los tres modelos SAM2 superan a
los dos de SAM v1, incluso cuando estos incorporan preentrenamiento
médico. SAM2.1-tiny \emph{general} ($0{,}9517$) supera a MedSAM v1
\emph{médico} ($0{,}9458$) por $+0{,}59$\,pp
($p = 2{,}5\times10^{-42}$). Segundo, el preentrenamiento médico
ayuda, pero solo dentro de la misma generación: MedSAM2-tiny supera
a SAM2.1-tiny por $+0{,}39$\,pp ($p = 9{,}5\times10^{-23}$),
mientras que en SAM v1 no compensa la arquitectura más antigua.
Tercero, el tamaño de \emph{backbone} no aporta nada en este régimen
\emph{box-prompted}: SAM2.1-tiny iguala o supera a SAM2.1-Large
($+0{,}14$\,pp, $p = 1{,}6\times10^{-7}$) con $5{,}7$ veces menos
parámetros y $3{,}7$ veces menos latencia. La mejor configuración,
con el Dice más alto y casi la menor latencia, es MedSAM2-tiny
($0{,}9556$, $26$\,ms): combina las dos palancas favorables
---generación SAM2 y preentrenamiento médico--- en el
\emph{backbone} más compacto. Las magnitudes absolutas son moderadas
por el techo alto del régimen \emph{box-prompted} con \emph{bounding
box} de referencia, pero todos los contrastes pareados son
significativos sobre las $1\,000$ imágenes del test. Que la
generación de arquitectura pese más que el preentrenamiento médico
matiza la hipótesis central del trabajo: los resultados sugieren
que, en este régimen \emph{box-prompted}, el diseño del segmentador
domina sobre la especialización dermatológica. Antes de derivar una recomendación
operativa de esta comparativa conviene acotar tres aspectos: cuánto
del resultado depende del \emph{prompt} y no del modelo, hasta qué
punto la cifra es estable, y dónde se sitúa su techo real.

\subsection{Valor del prompt: control automático sin caja}
\label{sec:seg-control-automatico}

Toda la comparativa anterior opera en régimen \emph{box-prompted}:
cada predicción parte de una \emph{bounding box} derivada de la
máscara de referencia. Para aislar la contribución de ese
\emph{prompt} se entrena un control automático que no recibe caja
alguna ---una U-Net con codificador ResNet-101 preentrenado en
ImageNet, ajuste fino denso de $51{,}5$\,M de parámetros sobre el
mismo dataset $v1c$--- y se evalúa sobre el test canónico de
ISIC2018. La tabla~\ref{tab:seg-automatico} contrasta su Dice con el
del mejor modelo \emph{promptable}.

\begin{table}[H]
\centering
\footnotesize
\setlength{\tabcolsep}{4pt}
\renewcommand{\arraystretch}{0.95}
\caption{Control automático \emph{frente a} régimen
\emph{promptable} sobre el test canónico de ISIC2018
($N = 1\,000$). La diferencia de Dice cuantifica el aporte
de la caja-\emph{prompt}.}
\label{tab:seg-automatico}
\begin{tabular}{llrrrr}
\hline
Modelo & Régimen & Dice & IoU & Precisión & Recall \\
\hline
MedSAM2-tiny      & \emph{promptable} & $\mathbf{0{,}9556}$ & $0{,}916$ & --- & --- \\
U-Net (ResNet-101) & automático        & $0{,}8811$ & $0{,}804$ & $0{,}856$ & $0{,}942$ \\
\hline
\end{tabular}
\end{table}

El control automático alcanza Dice $0{,}8811$, $7{,}5$\,pp por debajo
del mejor modelo \emph{box-prompted} ($0{,}9556$). La magnitud es
reveladora si se contrasta con la dispersión \emph{intra}-régimen: las
cinco arquitecturas \emph{promptables} de la
tabla~\ref{tab:seg-arquitecturas} caben en $\sim 1{,}0$\,pp de Dice,
mientras que retirar el \emph{prompt} cuesta entre seis y siete veces
más. El factor dominante del rendimiento no es, por tanto, la
arquitectura ni el preentrenamiento del segmentador, sino la
disponibilidad de la caja. El perfil de error del modo automático lo
confirma: recall alto ($0{,}942$) pero precisión que se desploma
($0{,}856$) y una distancia de Hausdorff~95 entre $2{,}5$ y $3$ veces
peor, con desviación estándar $0{,}121$ ---señal de fallos
catastróficos puntuales que la caja elimina al acotar la región de
interés. El control valida además que la comparativa $F2$--$F5$ fue un
contraste justo: todos los modelos compartían exactamente el mismo
\emph{prompt}, de modo que sus diferencias reflejan el segmentador y
no la información de localización.

\subsection{Robustez del resultado: validación cruzada, prompt y calibración}
\label{sec:seg-robustez}

La cifra de $0{,}9556$ procede de una única partición. Para acotar su
incertidumbre se somete el mejor modelo ---MedSAM2-tiny--- a una
validación cruzada de cinco pliegues estratificados por dataset de
origen, reentrenando el decodificador en cada pliegue y evaluando
siempre sobre el mismo test canónico blindado. La media sobre los
cinco pliegues es Dice $0{,}9554 \pm 0{,}0010$, con intervalo de
confianza del $95\%$ $[0{,}9543,\, 0{,}9564]$; la cifra principal
$0{,}9556$ cae dentro del intervalo, y la desviación
estándar inter-pliegue ($8\times 10^{-4}$) descarta que la cifra
dependa de una partición afortunada. Una ablación paralela del rango
de \emph{LoRA} ($r\in\{4,8,16,32\}$) sobre el mismo modelo arroja una
amplitud de $0{,}24$\,pp sin tendencia monótona ---por debajo del
ruido inter-pliegue---: el ajuste de bajo rango iguala en exactitud al
ajuste denso del decodificador sin ventaja de latencia ni un ahorro
de parámetros relevante.

La estabilidad frente a la calidad del \emph{prompt} es crítica para
el despliegue, donde la caja no procede de la máscara de referencia
sino de un anotador o de un detector automático. La
figura~\ref{fig:seg-box-robustness} reporta el Dice del modelo bajo
perturbaciones controladas de la caja. El desplazamiento aleatorio de
esquinas (\emph{jitter} de $\pm 20$ píxeles) apenas resta
$0{,}86$\,pp: la posición exacta de la esquina no es crítica. La
escala de la caja, en cambio, sí lo es, y de forma asimétrica.
Apretar la caja resulta catastrófico ---$-6{,}8$\,pp al $-10\%$ y
$-15{,}4$\,pp al $-20\%$, porque recorta lesión real--- mientras que
aflojarla degrada poco ---$-2{,}8$\,pp al $+10\%$, ya que el modelo
descarta el margen sobrante. Bajo una caja con holgura moderada
($+10\%$) el Dice se mantiene en torno a $0{,}93$. La lectura
operativa de esta asimetría se desarrolla en el
capítulo~\ref{cap:discusion} (sección~\ref{sec:disc-cv-robustez}).

\begin{figure}[H]
\centering
\includegraphics[width=0.82\textwidth]{fig_seg_box_robustness.png}
\caption{Robustez de MedSAM2-tiny a la calidad de la caja-\emph{prompt}
sobre el test canónico de ISIC2018 ($N = 1\,000$). El \emph{jitter}
de esquinas es prácticamente inocuo ($-0{,}86$\,pp), pero la escala
es crítica y asimétrica: apretar la caja degrada mucho más que
aflojarla.}
\label{fig:seg-box-robustness}
\end{figure}

El resto de comprobaciones de robustez confirman un modelo bien
comportado. La calibración a nivel de píxel es buena (error de
calibración esperado $0{,}0082$ frente a la máscara de referencia).
La equidad por fototipo, estimada sobre el ángulo tipológico
individual (ITA) como aproximación a la escala Fitzpatrick, muestra
una brecha de solo $1{,}1$\,pp entre el grupo más y menos favorecido,
con un descenso leve en los fototipos más oscuros. El aumento en test
(TTA) y el \emph{ensemble} de los cinco pliegues aportan mejoras
marginales ($+0{,}13$ y $+0{,}14$\,pp), que no compensan su coste, de
modo que la recomendación se mantiene en un único modelo.

\subsection{El techo de la métrica: acuerdo inter-anotador}
\label{sec:seg-techo-anotador}

Una cifra de Dice se mide siempre contra una máscara de referencia
trazada por una persona, y el contorno de una lesión es intrínsecamente
ambiguo: distintos especialistas dibujan bordes distintos. Para situar
el resultado frente a ese límite se emplea el subconjunto de IMA++ con
anotación múltiple ($312$ imágenes del test con dos o más máscaras
independientes), reservado y nunca usado en entrenamiento. El acuerdo
entre anotadores humanos ---Dice promedio de una máscara contra otra---
es de solo $0{,}728$. El modelo concuerda con el consenso de los
anotadores en $0{,}915$, y su predicción cae dentro de la envolvente de
las máscaras humanas el $97{,}8\%$ de las veces. Es decir, el modelo
coincide con el consenso \emph{mejor} de lo que los humanos coinciden
entre sí.

La lectura es que el $0{,}9556$ se sitúa esencialmente en el techo de
acuerdo que la propia tarea permite: cuando la referencia es la
conjetura difusa de un único anotador, el «error» residual es en buena
parte ruido de etiqueta, no deficiencia del modelo, y ningún modelo
alcanzaría $1{,}0$ contra una referencia ambigua. Las diferencias de
$\sim 1$\,pp entre arquitecturas
(tabla~\ref{tab:seg-arquitecturas}) son, de hecho, menores que el
desacuerdo inter-anotador, lo que recomienda no sobreinterpretarlas.
El matiz honesto es que el modelo es algo \emph{sobreconfiado} en los
bordes ambiguos: su error de calibración contra la máscara de consenso
\emph{blanda} sube a $0{,}094$, frente al $0{,}0082$ contra una
referencia binaria. En conjunto, estos resultados sugieren que la
principal limitación ya no reside en la capacidad del segmentador,
sino en la propia incertidumbre inherente a la anotación humana de
los bordes de lesión. Trazado este techo en una tarea con
referencia geométrica, queda por comprobar qué ocurre cuando la
supervisión llega en forma de lenguaje, lo que motiva el bloque
sobre clasificación \emph{zero-shot} multimodal.

\section{Clasificación zero-shot multimodal}
\label{sec:zs-results}

¿De qué depende que un modelo visión-lenguaje clasifique bien sin
entrenar cabeza alguna? Importa porque el \emph{zero-shot} promete
despliegue inmediato sin etiquetas. Se descubre que el factor
crítico no es la imagen sino la compatibilidad entre el tokenizador
del texto y el espacio del encoder. La tabla~\ref{tab:zs-ham} lo
cuantifica sobre HAM10000 ($7$ clases) con tres configuraciones de
\emph{prompts} y dos tokenizadores distintos.

\begin{table}[H]
\centering
\footnotesize
\setlength{\tabcolsep}{4pt}
\renewcommand{\arraystretch}{0.95}
\caption{Zero-shot multimodal sobre HAM10000 ($N_{\mathrm{test}} = 1\,232$,
$7$ clases). Mejor valor por columna en negrita.}
\label{tab:zs-ham}
\begin{tabular}{llllrr}
\hline
Modelo                    & Tokenizador  & \emph{Prompts}     & Acc       & AUROC     & F1-macro \\
\hline
DermLIP v1/v2             & PubMedBERT   & Genéricos          & $0{,}086$ & $0{,}366$ & $0{,}023$ \\
DermLIP v1/v2             & PubMedBERT   & Derm1M A3          & ---       & $0{,}516$ & --- \\
DermLIP original          & GPT-2/CLIP   & Derm1M A3          & $\mathbf{0{,}427}$ & $\mathbf{0{,}854}$ & --- \\
\hline
\end{tabular}
\end{table}

El rendimiento depende drásticamente de cómo se construye el texto.
Pasar de \emph{prompts} genéricos en inglés con tokenizador
PubMedBERT al conjunto A3 de Derm1M con tokenizador GPT-2/CLIP eleva
la AUROC de $0{,}366$ a $0{,}854$ ($+48{,}8$\,pp), como ilustra la
figura~\ref{fig:zs-prompts}. Los
\emph{prompts} genéricos con PubMedBERT quedan en AUROC$<0{,}5$,
señal de que proyectan en un espacio no alineado con el encoder
visual de la variante v1/v2. Cambiar solo el conjunto A3, sin tocar
el tokenizador, sube la AUROC a $0{,}516$; cambiar solo el
tokenizador, manteniendo A3, la lleva a $0{,}854$. El factor crítico
es, por tanto, la compatibilidad del tokenizador con el espacio de
\emph{embeddings}.

Este resultado tiene una implicación relevante para la
interpretación de los sistemas visión-lenguaje en dermatología. La
diferencia de casi $50$ puntos porcentuales de AUROC se obtiene sin
modificar el encoder visual ni las imágenes de entrada, lo que
indica que el principal cuello de botella no reside en la capacidad
de representación visual, sino en la alineación entre las
representaciones textuales y visuales. En consecuencia, evaluar un
modelo zero-shot únicamente a partir de una configuración de
\emph{prompts} o tokenizadores puede conducir a conclusiones
erróneas sobre sus capacidades reales.

Una vez identificado el papel determinante del tokenizador sobre
HAM10000, se evalúa hasta qué punto este comportamiento se mantiene
en escenarios clínicos externos. Para ello se considera la tarea
binaria melanoma/no-melanoma sobre cuatro datasets independientes
(DDI, MSKCC, HIBA y Derm7pt clínico) con la configuración
PubMedBERT de DermLIP v1/v2 (tabla~\ref{tab:zs-binaria}).

\begin{table}[H]
\centering
\footnotesize
\setlength{\tabcolsep}{4pt}
\caption{Zero-shot binario melanoma/no-melanoma con DermLIP v1/v2
(\emph{prompts} genéricos).}
\label{tab:zs-binaria}
\begin{tabular}{lrrrrr}
\hline
Dataset           & $N_{\mathrm{test}}$ & Acc       & BAcc      & AUROC     & F1-macro \\
\hline
DDI               & $137$               & $0{,}766$ & $0{,}499$ & $0{,}436$ & $0{,}463$ \\
MSKCC             & $1\,664$            & $0{,}716$ & $0{,}500$ & $0{,}430$ & $0{,}417$ \\
HIBA              & $334$               & $0{,}859$ & $0{,}497$ & $0{,}502$ & $0{,}462$ \\
Derm7pt clín.     & $168$               & $0{,}589$ & $0{,}472$ & $0{,}337$ & $0{,}439$ \\
\hline
\end{tabular}
\end{table}

En la formulación binaria, las cifras de BAcc próximas a
$0{,}500$ sobre los cuatro datasets indican que la rama contrastiva
con \emph{prompts} genéricos opera al nivel del azar. En conjunto,
los resultados muestran que el rendimiento zero-shot de DermLIP
depende mucho más de la formulación textual que de la capacidad
visual del encoder: la sensibilidad extrema al tokenizador y a los
\emph{prompts} limita la robustez del paradigma zero-shot puro y
sugiere que la transferencia multimodal requiere mecanismos de
alineamiento más sólidos para generalizar de forma fiable fuera del
dominio de entrenamiento. El detalle del análisis extendido sobre
sensibilidad a \emph{prompts} y tokenizador se desarrolla en el
Anexo~\ref{cap:anexog}.

\begin{figure}[H]
\centering
\includegraphics[width=0.78\textwidth]{fig_zs_prompts.png}
\caption{Sensibilidad zero-shot al tokenizador y a la
formulación del \emph{prompt} sobre HAM10000. La AUROC oscila
entre $0{,}366$ (DermLIP v1/v2 + PubMedBERT + prompts
genéricos) y $0{,}854$ (DermLIP original +
GPT-2/CLIP + prompts A3 de Derm1M), un rango de
$48{,}8$\,pp para variaciones que no afectan al encoder
visual.}
\label{fig:zs-prompts}
\end{figure}

\section{Comparación con modelos de lenguaje multimodales}
\label{sec:llm-results}

El bloque anterior mostró que la rama contrastiva de DermLIP es muy
sensible al texto. Cabe preguntarse si los grandes modelos de
lenguaje multimodales, entrenados sobre datasets mucho mayores,
cierran esa brecha frente a los especialistas. La
tabla~\ref{tab:llm-comparativa} lo contrasta evaluando seis modelos
de lenguaje multimodales sobre HAM10000, PAD-UFES-20 y DDI, junto a
los modelos de referencia de sección~\ref{sec:rep-panderm} y
sección~\ref{sec:ft-results}.

\begin{table}[H]
\centering
\footnotesize
\setlength{\tabcolsep}{4pt}
\renewcommand{\arraystretch}{0.95}
\caption{Comparación de paradigmas sobre HAM10000, PAD-UFES-20 y DDI
($N_{\mathrm{test}} = 1\,232 / 461 / 137$). FT: ajuste fino, LP:
sondeo lineal, ZS: zero-shot.}
\label{tab:llm-comparativa}
\begin{tabular}{llrrrrrr}
\hline
Modelo                 & Modo  & HAM Acc       & HAM BAcc      & PAD Acc       & PAD BAcc      & DDI Acc       & DDI BAcc \\
\hline
PanDerm Base (TTA HAM) & FT    & $\mathbf{0{,}920}$ & $\mathbf{0{,}852}$ & $0{,}755$     & $\mathbf{0{,}704}$ & $0{,}774$     & $0{,}693$ \\
PanDerm Large          & LP    & $0{,}888$     & $0{,}575$     & $\mathbf{0{,}772}$ & $0{,}642$     & $\mathbf{0{,}847}$ & $\mathbf{0{,}764}$ \\
SigLIP-Large           & LP    & $0{,}900$     & ---           & $0{,}718$     & ---           & $0{,}789$     & --- \\
MedGemma 27B + LoRA    & FT    & $0{,}802$     & $0{,}270$     & $0{,}430$     & $0{,}347$     & $0{,}657$     & $0{,}618$ \\
MedGemma 27B           & ZS    & $0{,}665$     & $0{,}243$     & $0{,}447$     & $0{,}342$     & $0{,}474$     & $0{,}566$ \\
GPT-4o                 & ZS    & $0{,}485$     & $0{,}465$     & $0{,}553$     & $0{,}523$     & $0{,}723$     & $0{,}648$ \\
GPT-4o-mini            & ZS    & $0{,}256$     & $0{,}319$     & $0{,}390$     & $0{,}435$     & $0{,}737$     & $0{,}556$ \\
Gemini 2.5 Pro         & ZS    & $0{,}403$     & ---           & $0{,}477$     & ---           & $0{,}441$     & --- \\
Gemini 2.5 Flash       & ZS    & $0{,}380$     & ---           & $0{,}486$     & ---           & $0{,}555$     & --- \\
DermLIP                & ZS    & $0{,}239$     & $0{,}581$     & $0{,}601$     & ---           & $0{,}329$     & --- \\
BLIP-2 (Flan-T5-XL)    & ZS    & $0{,}015$     & $0{,}145$     & $0{,}017$     & $0{,}167$     & $0{,}796$     & $0{,}542$ \\
InstructBLIP (Flan-T5-XL)& ZS  & $0{,}024$     & $0{,}123$     & $0{,}171$     & $0{,}224$     & $0{,}788$     & $0{,}500$ \\
\hline
\end{tabular}
\end{table}

Los resultados muestran una separación clara entre los modelos
especializados y los modelos multimodales generalistas. El mejor
sistema dermatológico (PanDerm Base ajustado con TTA) alcanza una
accuracy de $0{,}920$ sobre HAM10000, mientras que el mejor modelo
multimodal generalista en régimen \emph{zero-shot} (GPT-4o) obtiene
$0{,}485$: una diferencia de $43{,}5$\,pp que se mantiene de forma
consistente en los datasets evaluados. La
figura~\ref{fig:llm-comparison} sintetiza este hallazgo principal.
En el extremo inferior, BLIP-2 e InstructBLIP descienden hasta
accuracy $\sim 0{,}02$ sobre HAM10000.

El resultado sugiere que el conocimiento visual especializado
adquirido mediante preentrenamiento dermatológico continúa siendo
más relevante para la clasificación de lesiones cutáneas que el
conocimiento generalista acumulado por los grandes modelos
multimodales. Aunque estos últimos poseen capacidades lingüísticas y
de razonamiento considerablemente superiores, dicha ventaja no
compensa la ausencia de una adaptación específica al dominio
dermatológico.

\subsection{Ajuste supervisado de un LLM multimodal: MedGemma con LoRA}

¿Reduce la brecha adaptar el LLM al dominio? La adaptación con LoRA
sobre MedGemma~27B, usando el \emph{split} de entrenamiento de
HAM10000, mejora $+13{,}7$\,pp de accuracy sobre HAM10000 y
$+18{,}3$\,pp sobre DDI respecto a su variante \emph{zero-shot},
pero no aporta sobre PAD-UFES-20 ($-1{,}7$\,pp). El ajuste
supervisado reduce, por tanto, parcialmente la distancia, pero no
alcanza a los encoders especializados: incluso con LoRA, MedGemma~27B
($0{,}802$ de accuracy sobre HAM10000) queda $11{,}8$\,pp por debajo
del mejor especialista.

El desglose de errores por clase de GPT-4o ---con patrones
asimétricos como \emph{recall} $0{,}000$ sobre \emph{actinic
keratosis} y un fuerte sesgo de predicción hacia categorías menos
restringidas en el \emph{prompt}--- y el coste de inferencia de cada
modelo (APIs comerciales y ejecución local) se detallan en el
Anexo~\ref{cap:anexog}.

\begin{figure}[H]
\centering
\includegraphics[width=0.92\textwidth]{fig_llm_comparison.png}
\caption{Comparativa de paradigmas sobre HAM10000. Gradiente
de $43{,}5$\,pp accuracy entre el mejor especialista
ajustado (PanDerm Base FT con TTA, $0{,}920$) y el mejor LLM
multimodal generalista \emph{zero-shot} (GPT-4o,
$0{,}485$). BLIP-2 e InstructBLIP descienden hasta accuracy
$\sim 0{,}02$.}
\label{fig:llm-comparison}
\end{figure}

\section{Equidad por fototipo cutáneo}
\label{sec:equidad-results}

¿Rinde un modelo de forma equitativa entre fototipos cutáneos, y
cómo se distingue una equidad real de una métrica engañosa? La
cuestión es ética y clínica a la vez. El análisis revela que un
\emph{gap} pequeño puede esconder un modelo uniformemente peor, de
modo que la métrica aislada no decide; la comparación de cinco
encoders y el test de DeLong lo establecen. La
tabla~\ref{tab:equidad-global} reporta el rendimiento global de
cinco encoders sobre Fitzpatrick17k completo ($N = 16\,577$) en
las dos formulaciones consideradas: clasificación a $114$
patologías (grano fino del dataset) y clasificación a $3$ clases
de malignidad (benigno / no~neoplásico / maligno).

\begin{table}[H]
\centering
\footnotesize
\setlength{\tabcolsep}{4pt}
\renewcommand{\arraystretch}{0.95}
\caption{Equidad sobre Fitzpatrick17k completo. Rendimiento global
en las dos formulaciones. Mejor valor por columna en negrita.}
\label{tab:equidad-global}
\begin{tabular}{lrrrrrrrr}
\hline
                          & \multicolumn{4}{c}{$114$ patologías} & \multicolumn{4}{c}{$3$ clases malignidad} \\
\cline{2-5}\cline{6-9}
Modelo                    & Acc           & BAcc          & AUROC         & W-F1          & Acc           & BAcc          & AUROC         & W-F1 \\
\hline
PanDerm Large             & $\mathbf{0{,}581}$ & $0{,}549$     & $0{,}968$     & $\mathbf{0{,}576}$ & $0{,}835$     & $0{,}699$     & $0{,}908$     & $0{,}828$ \\
PanDerm Base              & $0{,}500$     & $0{,}478$     & $0{,}955$     & $0{,}491$     & $0{,}819$     & $0{,}657$     & $0{,}882$     & $0{,}807$ \\
DermLIP v2 (visual)       & $0{,}572$     & $\mathbf{0{,}549}$ & $\mathbf{0{,}971}$ & $0{,}569$ & $\mathbf{0{,}852}$ & $\mathbf{0{,}721}$ & $\mathbf{0{,}909}$ & $\mathbf{0{,}846}$ \\
DINOv2 ViT-L/14           & $0{,}519$     & $0{,}485$     & $0{,}947$     & $0{,}516$     & $0{,}809$     & $0{,}656$     & $0{,}854$     & $0{,}801$ \\
BiomedCLIP                & $0{,}334$     & $0{,}321$     & $0{,}910$     & $0{,}323$     & $0{,}799$     & $0{,}590$     & $0{,}830$     & $0{,}776$ \\
\hline
\end{tabular}
\end{table}

La tabla~\ref{tab:equidad-fototipo} reporta la BAcc por fototipo
Fitzpatrick sobre la formulación de tres clases.

\begin{table}[H]
\centering
\footnotesize
\setlength{\tabcolsep}{4pt}
\renewcommand{\arraystretch}{0.95}
\caption{BAcc por subgrupo Fitzpatrick (3 clases de malignidad) y
\emph{gap} entre fototipos extremos. Tamaños muestrales por
subgrupo: $299 / 470 / 325 / 261 / 169 / 73$ para FST~I a~VI.}
\label{tab:equidad-fototipo}
\begin{tabular}{lrrrrrrr}
\hline
Modelo                    & FST I         & FST II        & FST III       & FST IV        & FST V         & FST VI        & Gap I-VI \\
\hline
PanDerm Large             & $0{,}723$     & $0{,}708$     & $0{,}743$     & $0{,}669$     & $0{,}568$     & $0{,}506$     & $+0{,}217$ \\
PanDerm Base              & $0{,}653$     & $0{,}693$     & $0{,}669$     & $0{,}630$     & $0{,}585$     & $0{,}384$     & $+0{,}269$ \\
DermLIP v2 (visual)       & $0{,}716$     & $0{,}730$     & $0{,}765$     & $0{,}710$     & $0{,}693$     & $0{,}434$     & $+0{,}281$ \\
DINOv2 ViT-L/14           & $0{,}673$     & $0{,}671$     & $0{,}679$     & $0{,}654$     & $0{,}553$     & $0{,}368$     & $+0{,}306$ \\
BiomedCLIP                & $0{,}569$     & $0{,}602$     & $0{,}606$     & $0{,}566$     & $0{,}590$     & $\mathbf{0{,}445}$ & $\mathbf{+0{,}124}$ \\
\hline
\end{tabular}
\end{table}

Los cinco modelos tienen \emph{gap}~I-VI positivo, es decir, rinden
peor sobre fototipos elevados. DermLIP v2 logra la mejor BAcc global
($0{,}721$ sobre $3$ clases) y la mejor AUROC global ($0{,}909$),
pero su \emph{gap}~I-VI ($+0{,}281$) supera al de PanDerm Large
($+0{,}217$). BiomedCLIP tiene el \emph{gap} menor ($+0{,}124$),
pero a costa de una BAcc global inferior ($0{,}590$). Este contraste
muestra que el \emph{gap} aislado no decide nada: un modelo
uniformemente peor entre subgrupos puede exhibir un \emph{gap} bajo
sin ser equitativo en sentido sustantivo
(figura~\ref{fig:fairness-5models}). La AUROC global de PanDerm
Large sobre $3$ clases es $0{,}908$, con un \emph{gap} de AUROC
entre fototipos extremos de $0{,}089$ (FST~IV: $0{,}919$ vs FST~VI:
$0{,}830$).

La tabla~\ref{tab:cross-eval} amplía la evaluación con un
experimento cruzado \emph{Light}~$\leftrightarrow$~\emph{Dark}, donde
\emph{Light} agrupa FST~I-III y \emph{Dark} agrupa FST~IV-VI.

\begin{table}[H]
\centering
\footnotesize
\setlength{\tabcolsep}{4pt}
\caption{Evaluación cruzada \emph{Light}~$\leftrightarrow$~\emph{Dark}
(BAcc, formulación de $3$ clases). \emph{L}: \emph{Light};
\emph{D}: \emph{Dark}; \emph{All}: ambos grupos.}
\label{tab:cross-eval}
\begin{tabular}{lrrrrrr}
\hline
Modelo                    & L$\to$L       & L$\to$D       & D$\to$D       & D$\to$L       & All$\to$L     & All$\to$D \\
\hline
PanDerm Large             & $0{,}699$     & $0{,}589$     & $0{,}701$     & $0{,}667$     & $0{,}726$     & $0{,}627$ \\
PanDerm Base              & $0{,}673$     & $0{,}557$     & $0{,}621$     & $0{,}629$     & $0{,}673$     & $0{,}594$ \\
DermLIP v2 (visual)       & $\mathbf{0{,}729}$ & $\mathbf{0{,}651}$ & $0{,}667$ & $\mathbf{0{,}683}$ & $\mathbf{0{,}745}$ & $\mathbf{0{,}682}$ \\
DINOv2 ViT-L/14           & $0{,}655$     & $0{,}523$     & $0{,}625$     & $0{,}569$     & $0{,}682$     & $0{,}589$ \\
BiomedCLIP                & $0{,}595$     & $0{,}532$     & $0{,}571$     & $0{,}551$     & $0{,}594$     & $0{,}565$ \\
\hline
\end{tabular}
\end{table}

Entrenar solo con fototipos claros y evaluar sobre oscuros (L$\to$D
frente a L$\to$L) penaliza el rendimiento. La caída es mayor en
DINOv2 ($\Delta = -0{,}132$\,pp BAcc) y menor en BiomedCLIP
($\Delta = -0{,}063$\,pp). DermLIP v2 muestra la menor caída
Dark$\to$Light frente a Dark$\to$Dark ($\Delta = +0{,}016$\,pp), lo
que apunta a representaciones más invariantes al fototipo cutáneo
(figura~\ref{fig:cross-eval}).

\begin{figure}[H]
\centering
\includegraphics[width=0.95\textwidth]{figura_fairness_5models.png}
\caption{Equidad por fototipo Fitzpatrick (formulación de tres
clases de malignidad) sobre los cinco encoders evaluados. Los
cinco presentan \emph{gap}~I-VI positivo (entre $+0{,}124$ y
$+0{,}306$\,pp BAcc). Paradoja BiomedCLIP: \emph{gap} mínimo
asociado a peor BAcc global, lo que demuestra el carácter no
decisional de la métrica aislada.}
\label{fig:fairness-5models}
\end{figure}

\begin{figure}[H]
\centering
\includegraphics[width=0.85\textwidth]{fig_cross_eval.png}
\caption{Evaluación cruzada
\emph{Light}~$\leftrightarrow$~\emph{Dark} sobre Fitzpatrick17k
(BAcc en formulación de tres clases). La caída
L$\to$D vs L$\to$L cuantifica la pérdida de rendimiento al
entrenar sólo con fototipos claros y evaluar sobre fototipos
oscuros, máxima en DINOv2 ($\Delta = -0{,}132$\,pp).}
\label{fig:cross-eval}
\end{figure}

Las diferencias observadas entre encoders en las tablas anteriores
permiten establecer un \emph{ranking} descriptivo de rendimiento y
equidad. Sin embargo, la proximidad entre algunos resultados
---especialmente entre PanDerm Large y DermLIP v2--- plantea una
cuestión inmediata: ¿reflejan estas diferencias una ventaja real o
son compatibles con la variabilidad muestral del conjunto de
evaluación? Para responder a esta pregunta se realiza un contraste
formal mediante el test de DeLong.

\subsection{Validación estadística del ranking de encoders: test de
DeLong}
\label{sec:equidad-delong}

Las diferencias de AUROC entre encoders de las tablas anteriores se
contrastan formalmente sobre un mismo conjunto de test mediante el
test de DeLong~\cite{delong1988}, que controla la correlación
inducida por evaluar los cinco modelos sobre las mismas imágenes. El
contraste se realiza sobre el endpoint binario de detección de
malignidad (clase \emph{maligno} frente al resto) en el conjunto de
test de Fitzpatrick17k ($N = 1\,658$, $226$ casos malignos),
entrenando para cada encoder una regresión logística binaria sobre
sus \emph{features} congeladas con el protocolo de
sección~\ref{sec:lp}. La tabla~\ref{tab:delong-auroc} reporta la
AUROC binaria y su intervalo de confianza al $95\%$ derivado de la
varianza de DeLong; la tabla~\ref{tab:delong-pares} reporta los
$p$-valores pareados (dos colas).

\begin{table}[H]
\centering
\footnotesize
\setlength{\tabcolsep}{4pt}
\renewcommand{\arraystretch}{0.95}
\caption{AUROC binaria de detección de malignidad (\emph{maligno}
frente al resto) sobre el test de Fitzpatrick17k
($N = 1\,658$, $226$ malignos), con intervalo de confianza al
$95\%$ por la varianza de DeLong. Mejor valor en negrita.}
\label{tab:delong-auroc}
\begin{tabular}{lrl}
\hline
Modelo            & AUROC               & IC95\% (DeLong) \\
\hline
PanDerm Large     & $\mathbf{0{,}9488}$ & $[0{,}9350,\, 0{,}9626]$ \\
DermLIP v2        & $0{,}9452$          & $[0{,}9307,\, 0{,}9598]$ \\
PanDerm Base      & $0{,}9296$          & $[0{,}9130,\, 0{,}9463]$ \\
DINOv2 ViT-L/14   & $0{,}9150$          & $[0{,}8932,\, 0{,}9367]$ \\
BiomedCLIP        & $0{,}8915$          & $[0{,}8672,\, 0{,}9158]$ \\
\hline
\end{tabular}
\end{table}

\begin{table}[H]
\centering
\footnotesize
\setlength{\tabcolsep}{4pt}
\renewcommand{\arraystretch}{0.95}
\caption{Test de DeLong pareado entre encoders sobre la AUROC binaria
de malignidad: $p$-valor a dos colas. Valores no significativos al
nivel nominal $0{,}05$ en cursiva. \emph{Provenance}:
\texttt{delong\_malignancy\_results.json} del repositorio de
investigación.}
\label{tab:delong-pares}
\begin{tabular}{lrrrr}
\hline
                & DermLIP v2          & PanDerm Base        & DINOv2              & BiomedCLIP \\
\hline
PanDerm Large   & \emph{$0{,}55$}     & $1{,}2\times 10^{-3}$ & $4{,}2\times 10^{-5}$ & $8{,}7\times 10^{-9}$ \\
DermLIP v2      & ---                 & $7{,}3\times 10^{-3}$ & $2{,}4\times 10^{-3}$ & $6{,}0\times 10^{-7}$ \\
PanDerm Base    & ---                 & ---                 & \emph{$0{,}12$}     & $1{,}2\times 10^{-4}$ \\
DINOv2          & ---                 & ---                 & ---                 & \emph{$0{,}066$} \\
\hline
\end{tabular}
\end{table}

Tres lecturas se desprenden del contraste. \textit{Primero}, los dos
encoders específicos de dermatología, PanDerm Large y DermLIP v2,
empatan estadísticamente ($\Delta\mathrm{AUROC} = 0{,}0036$,
$p = 0{,}55$): ningún dato sostiene una jerarquía entre ambos sobre
este endpoint. \textit{Segundo}, ambos superan de forma
significativa a los codificadores generalistas DINOv2 y BiomedCLIP
y a PanDerm Base ($p \leq 7{,}3\times 10^{-3}$ en los seis
contrastes), lo que respalda cuantitativamente la ventaja del
preentrenamiento específico de dominio. \textit{Tercero}, bajo
corrección conservadora de Bonferroni para las diez comparaciones
(umbral $0{,}005$), todos los contrastes de los dos líderes frente
a los generalistas se mantienen significativos salvo DermLIP~v2
frente a PanDerm~Base ($p = 7{,}3\times 10^{-3}$, significativo solo
a nivel nominal); las diferencias DINOv2--BiomedCLIP y
DINOv2--PanDerm~Base no son significativas, lo que delimita qué
distinciones del \emph{ranking} son atribuibles a señal y cuáles a
ruido muestral.

En síntesis, el contraste establece dos conclusiones con respaldo
estadístico: entre los dos encoders específicos de dermatología no
existe un ganador ---empatan en $p = 0{,}55$, y ningún dato sostiene
una jerarquía entre ellos---, mientras que ambos superan de forma
significativa a los codificadores generalistas. Este segundo
resultado confirma, ahora mediante un contraste formal controlado
por la correlación entre modelos, la ventaja del preentrenamiento
específico de dominio que el capítulo documenta de manera recurrente
en transferencia, clasificación y razonamiento multimodal.

\section{Validación en castellano sobre DermapixelAI 1.0}
\label{sec:dermapixel-results}

Los bloques anteriores evalúan los modelos sobre datasets externos
anglosajones. El objetivo de generalización planteado en la
sección~\ref{sec:pregunta} exige un contraste adicional: comprobar
si los hallazgos se sostienen sobre material clínico en castellano,
construido de forma independiente y libre de fuga frente al
preentrenamiento. DermapixelAI~1.0 reúne ambas condiciones y motiva
el cambio de objeto de esta sección, del benchmark externo al
dataset propio.

El conjunto de experimentos sobre DermapixelAI~1.0 que abre esta
sección y se extiende hasta la sección~\ref{sec:spanderm-clip}
responde a una única pregunta ---si los modelos fundacionales
dermatológicos transfieren al material clínico en castellano--- y
comparte un hallazgo transversal que conviene anticipar: no existe
un método óptimo único, sino que la mejor configuración depende del
nivel ontológico evaluado. La tabla~\ref{tab:dermapixel-taxonomia}
fija la nomenclatura de los cuatro sistemas contrastados a lo largo
del bloque para evitar ambigüedad entre ellos.

\begin{table}[H]
\centering
\footnotesize
\setlength{\tabcolsep}{4pt}
\renewcommand{\arraystretch}{0.95}
\caption{Nomenclatura de los cuatro sistemas evaluados sobre
DermapixelAI~1.0 a lo largo del bloque
secciones~\ref{sec:dermapixel-results}--\ref{sec:spanderm-clip}.}
\label{tab:dermapixel-taxonomia}
\begin{tabular}{lll}
\hline
Sistema & Qué es & Sección \\
\hline
PanDerm LP                          & Codificador congelado $+$ regresión logística           & §\ref{sec:dermapixel-results} \\
DermLIP \emph{zero-shot}            & Modelo visión-lenguaje externo (sin cabeza entrenada)   & §\ref{sec:dermapixel-zs} \\
Dermapixel R0 (cabeza superv.\ L2) & Rama supervisada con adaptación LoRA $r=16$             & §\ref{sec:dermapixel-spanderm-v0} \\
Dermapixel R0 (rama contrastiva)   & Rama \emph{CLIP-style} castellana \emph{zero-shot open} & §\ref{sec:spanderm-clip} \\
\hline
\end{tabular}
\end{table}

La auditoría anterior dejó a Dermnet como único dataset externo con
solapamiento exacto. DermapixelAI~1.0 ofrece, frente a ello, un
escenario limpio en castellano: a diferencia de Dermnet, con
solapamiento documentado frente a Derm1M
(sección~\ref{sec:auditoria-resultados}), el dataset DermapixelAI~1.0
descrito en el Anexo~\ref{cap:anexoe} es completamente independiente
del dataset de preentrenamiento de la familia PanDerm/DermLIP. La
auditoría por hash MD5 sobre $507\,657$ imágenes catalogadas
(Derm1M $415\,451$ ítems, HAM10000, BCN20000, ISIC2017/2018,
PH2, PAD-UFES-20, DDI, Dermnet, Fitzpatrick17k, HIBA, MSKCC y
Derm7pt) reporta intersección nula: $0$ coincidencias sobre las
$1\,062$ imágenes evaluables del dataset. La evaluación reportada en
esta sección no está, por tanto, sujeta a la fuga de información
caracterizada en sección~\ref{sec:auditoria-resultados}.

\subsection{Transferencia congelada: sondeo lineal y validación cruzada}

\paragraph{Sondeo lineal sobre el \emph{split} fijo.}
Se aplica el protocolo de sondeo lineal descrito en
sección~\ref{sec:lp} sobre las $1\,062$ imágenes con
\texttt{label\_source=ontology}, particionadas en train
($N = 874$), validación ($N = 152$) y test ($N = 36$). PanDerm
Base ($768$-D) y PanDerm Large ($1\,024$-D) operan como extractores
congelados; la cabeza es un clasificador \emph{LogisticRegression}
(solver L-BFGS, $C = 1{,}0$). La evaluación L2 se reporta sobre las
$38$ subcategorías efectivas tras consolidar las dos variantes
ortográficas de \emph{Trastornos de la queratinización}
(Anexo~\ref{cap:anexoc}, sección~\ref{sec:anexo-l2-l3}). Dos
patrones estructuran el resultado (tabla~\ref{tab:dermapixel-lp}).
Primero, PanDerm Large supera sistemáticamente a Base en los tres
niveles (entre $+7{,}8$ y $+9{,}9$\,pp de AUROC y entre $+4{,}4$ y
$+16{,}5$\,pp de BAcc). Segundo, el rendimiento disminuye de forma
monótona al aumentar la granularidad diagnóstica: la BAcc de
PanDerm Large cae de $0{,}735$ (L1, cuatro categorías etiológicas)
a $0{,}265$ (L2, $38$ subcategorías) y a $0{,}184$ (L3, $224$
diagnósticos específicos).

\begin{table}[H]
\centering
\footnotesize
\setlength{\tabcolsep}{4pt}
\renewcommand{\arraystretch}{0.95}
\caption{Sondeo lineal con PanDerm Base y Large sobre
DermapixelAI~1.0 en los tres niveles de la ontología jerárquica
($N_{\mathrm{train}} = 874$, $N_{\mathrm{test}} = 36$ en L1/L2 y
$N_{\mathrm{test}} = 28$ en L3 tras filtrar clases ausentes en
train). Entre corchetes, IC95\% por \emph{bootstrap} estratificado
($1\,000$ remuestreos). Mejor valor por columna y nivel en
negrita.}
\label{tab:dermapixel-lp}
\adjustbox{max width=\textwidth}{%
\begin{tabular}{llrrrrr}
\hline
Nivel & Modelo & Acc & BAcc & AUROC & W-F1 & Kappa \\
\hline
L1 (4 clases)   & Base  & $0{,}639$ [$0{,}472$--$0{,}778$]
                        & $0{,}570$ [$0{,}371$--$0{,}751$]
                        & $0{,}697$ [$0{,}527$--$0{,}841$]
                        & $0{,}624$ [$0{,}457$--$0{,}760$]
                        & $0{,}296$ \\
L1 (4 clases)   & Large & $\mathbf{0{,}750}$ [$0{,}639$--$0{,}861$]
                        & $\mathbf{0{,}735}$ [$0{,}637$--$0{,}835$]
                        & $\mathbf{0{,}796}$ [$0{,}656$--$0{,}930$]
                        & $\mathbf{0{,}722}$ [$0{,}596$--$0{,}845$]
                        & $\mathbf{0{,}520}$ \\
L2 (38 clases)  & Base  & $0{,}222$ [$0{,}139$--$0{,}333$]
                        & $0{,}221$ [$0{,}147$--$0{,}294$]
                        & $0{,}707$ [$0{,}666$--$0{,}747$]
                        & $0{,}212$ [$0{,}092$--$0{,}302$]
                        & $0{,}170$ \\
L2 (38 clases)  & Large & $\mathbf{0{,}250}$ [$0{,}167$--$0{,}333$]
                        & $\mathbf{0{,}265}$ [$0{,}206$--$0{,}338$]
                        & $\mathbf{0{,}793}$ [$0{,}749$--$0{,}837$]
                        & $\mathbf{0{,}223}$ [$0{,}139$--$0{,}291$]
                        & $\mathbf{0{,}199}$ \\
L3 (224 clases) & Base  & $0{,}143$ [$0{,}107$--$0{,}179$]
                        & $0{,}132$ [$0{,}105$--$0{,}158$]
                        & $0{,}735$ [$0{,}693$--$0{,}779$]
                        & $0{,}107$ [$0{,}064$--$0{,}141$]
                        & $0{,}112$ \\
L3 (224 clases) & Large & $\mathbf{0{,}179}$ [$0{,}143$--$0{,}214$]
                        & $\mathbf{0{,}184}$ [$0{,}158$--$0{,}211$]
                        & $\mathbf{0{,}813}$ [$0{,}769$--$0{,}856$]
                        & $\mathbf{0{,}143}$ [$0{,}100$--$0{,}184$]
                        & $\mathbf{0{,}154}$ \\
\hline
\end{tabular}}
\end{table}

PanDerm Large supera a Base en las cinco métricas y los tres
niveles (tabla~\ref{tab:dermapixel-lp}). La cifra L1 BAcc $= 0{,}735$
(IC95\% [$0{,}637$--$0{,}835$]) se sitúa dentro del rango reportado
para PanDerm Large sobre los diez datasets externos
(tabla~\ref{tab:lp-panderm}), lo que ubica a DermapixelAI~1.0 como un
\emph{benchmark} exigente en castellano sin dependencia de marcas
dermatoscópicas estandarizadas. Frente al azar, PanDerm Large
multiplica el rendimiento por $2{,}9\times / 10{,}2\times /
40{,}9\times$ en L1/L2/L3.


\paragraph{Validación cruzada agrupada por caso.}
El test fijo $N = 36$ corresponde a un único \emph{split} publicado.
Para estimar el rendimiento esperado con menor dependencia de la
realización muestral, se aplica una validación cruzada de cinco
\emph{folds} agrupada por \texttt{case\_id}
(\emph{StratifiedGroupKFold} estratificado por L1 sobre las
$1\,062$ imágenes), de modo que ninguna imagen del mismo caso
aparezca simultáneamente en \emph{train} y \emph{test}. La
tabla~\ref{tab:dermapixel-cv} reporta la media y la desviación
estándar sobre los cinco \emph{folds}.

\begin{table}[H]
\centering
\footnotesize
\setlength{\tabcolsep}{4pt}
\renewcommand{\arraystretch}{0.95}
\caption{Validación cruzada agrupada por caso (cinco \emph{folds},
estratificado por L1 sobre \texttt{case\_id}) con PanDerm Large
sobre DermapixelAI~1.0. Cifras: media $\pm$ desviación estándar
sobre los cinco \emph{folds}. La cardinalidad L3 reportada aquí
(250 clases) corresponde a las clases efectivas en el dataset
completo; las restantes tablas del split fijo
(tabla~\ref{tab:dermapixel-lp}, tabla~\ref{tab:dermapixel-ensemble}
y las del Anexo~\ref{cap:anexoj}) usan 224 clases L3, las
únicas representadas en el split \texttt{train}.}
\label{tab:dermapixel-cv}
\begin{tabular}{lrrrrrr}
\hline
Nivel           & Acc@1                 & Acc@3                 & BAcc                  & AUROC                 & W-F1                  & Kappa \\
\hline
L1 (4 cls)      & $0{,}574 \pm 0{,}029$ & $0{,}994 \pm 0{,}004$ & $0{,}440 \pm 0{,}058$ & $0{,}753 \pm 0{,}054$ & $0{,}569 \pm 0{,}025$ & $0{,}314 \pm 0{,}035$ \\
L2 (38 cls)     & $0{,}227 \pm 0{,}015$ & $0{,}418 \pm 0{,}037$ & $0{,}168 \pm 0{,}020$ & $0{,}749 \pm 0{,}019$ & $0{,}221 \pm 0{,}022$ & $0{,}186 \pm 0{,}015$ \\
L3 (250 cls)    & $0{,}124 \pm 0{,}007$ & $0{,}232 \pm 0{,}028$ & $0{,}108 \pm 0{,}014$ & $0{,}800 \pm 0{,}018$ & $0{,}124 \pm 0{,}005$ & $0{,}110 \pm 0{,}005$ \\
\hline
\end{tabular}
\end{table}

La validación cruzada agrupada por caso constituye la estimación
defendible del rendimiento esperado de un sondeo lineal congelado
sobre DermapixelAI~1.0. La sobrestimación del test fijo respecto a
la media de los cinco \emph{folds} es de $+29{,}5$\,pp BAcc en L1,
$+9{,}7$\,pp en L2 y $+7{,}6$\,pp en L3, atribuible a la
distribución favorable del split publicado (test sin la categoría
\emph{Genodermatosis} y baja densidad de la cola larga L3): el
split fijo es, por tanto, una estimación optimista. A diferencia
de la BAcc, la AUROC se mantiene estable frente al cambio de
partición, con caídas inferiores a $5$\,pp en los tres niveles
($0{,}796 \to 0{,}753$ en L1, $0{,}793 \to 0{,}749$ en L2 y
$0{,}813 \to 0{,}800$ en L3). Este contraste ---accuracy y BAcc
dependientes del \emph{split}, AUROC robusta--- es relevante para
comparar configuraciones sobre datasets de tamaño moderado. Las
cifras del test fijo (tabla~\ref{tab:dermapixel-lp}) se conservan
como punto operativo del split publicado y como base de la
comparación con el experimento \emph{zero-shot} de la
sección~\ref{sec:dermapixel-zs}, que reutiliza el mismo test para
garantizar la trazabilidad cruzada.


\subsection{Validación del procedimiento de etiquetado}

Una posible amenaza a la validez de DermapixelAI~1.0 es que parte de
las etiquetas etiológicas procede de la ontología y no de una
revisión individual realizada por la dermatóloga. Para cuantificar
este efecto se compara el rendimiento por sondeo lineal sobre el
subconjunto completamente verificado (\texttt{rosa\_verified=True})
frente al dataset completo, en los tres niveles ontológicos. Las
diferencias observadas son pequeñas, inconsistentes en dirección y
compatibles con la variabilidad muestral estimada mediante
\emph{bootstrap}: todos los intervalos de confianza al $95\%$
muestran un amplio solapamiento y no se observa un patrón consistente
de deterioro asociado al uso de etiquetas derivadas. Esto sugiere
que la imputación etiológica sobre las $170$ imágenes no revisadas
individualmente no introduce una degradación apreciable del
rendimiento, lo que respalda el flujo de etiquetado descrito en el
Anexo~\ref{cap:anexoe}. Los resultados completos por nivel
ontológico figuran en el Anexo~\ref{cap:anexoj}
(tabla~\ref{tab:dermapixel-rosa}).

Diversas ablaciones metodológicas adicionales sobre los mismos
\emph{embeddings} congelados ---cabezas alternativas ($k$-NN, MLP,
regresión logística reponderada), función de pérdida y
\emph{Test-Time Augmentation}--- no modificaron las conclusiones
principales y se documentan íntegramente en el
Anexo~\ref{cap:anexoj}.


\subsection{Combinación de codificadores}
\label{sec:dermapixel-ensemble}

Se evalúa si las representaciones de distintos codificadores
contienen información complementaria. Se entrena una cabeza de
regresión logística independiente por codificador (PanDerm Base,
PanDerm Large y DermLIP v2; protocolo de sección~\ref{sec:lp}) sobre
los \emph{embeddings} congelados, y se contrastan frente a cinco
estrategias de combinación (promedio de las tres distribuciones,
promedio Large$+$DermLIP, máximo por clase y apilado con
meta-cabeza). La tabla~\ref{tab:dermapixel-ensemble} reporta los
resultados sobre los tres niveles ontológicos en el mismo test fijo.

\begin{table}[H]
\centering
\footnotesize
\setlength{\tabcolsep}{4pt}
\renewcommand{\arraystretch}{0.95}
\caption{Combinación de codificadores sobre DermapixelAI~1.0. LP
individual sobre PanDerm Base, PanDerm Large y DermLIP v2; tres
estrategias de combinación: promedio Large$+$DermLIP, máximo por
clase entre los tres, y apilado con meta-cabeza. Mejor valor por
columna y nivel en negrita.}
\label{tab:dermapixel-ensemble}
\begin{tabular}{llrrrr}
\hline
Nivel & Estrategia                         & Acc@1              & Acc@3              & BAcc               & AUROC \\
\hline
L1 (4 cls)   & LP PanDerm Base             & $0{,}639$          & ---                & $0{,}570$          & $0{,}697$ \\
L1 (4 cls)   & LP PanDerm Large            & $\mathbf{0{,}750}$ & $\mathbf{1{,}000}$ & $\mathbf{0{,}735}$ & $0{,}796$ \\
L1 (4 cls)   & LP DermLIP v2               & $0{,}694$          & $\mathbf{1{,}000}$ & $0{,}680$          & $0{,}789$ \\
L1 (4 cls)   & Avg L$+$D                   & $\mathbf{0{,}750}$ & $\mathbf{1{,}000}$ & $0{,}702$          & $\mathbf{0{,}817}$ \\
L1 (4 cls)   & Avg 3                       & $0{,}722$          & $\mathbf{1{,}000}$ & $0{,}694$          & $0{,}811$ \\
L1 (4 cls)   & Máx 3                       & $0{,}694$          & $\mathbf{1{,}000}$ & $0{,}652$          & $0{,}805$ \\
L1 (4 cls)   & Stack 3                     & $0{,}694$          & $\mathbf{1{,}000}$ & $0{,}703$          & $0{,}790$ \\
\hline
L2 (38 cls)  & LP PanDerm Base             & $0{,}222$          & ---                & $0{,}221$          & $0{,}707$ \\
L2 (38 cls)  & LP PanDerm Large            & $0{,}250$          & $0{,}500$          & $0{,}265$          & $0{,}793$ \\
L2 (38 cls)  & LP DermLIP v2               & $\mathbf{0{,}333}$ & $0{,}528$          & $\mathbf{0{,}279}$ & $\mathbf{0{,}860}$ \\
L2 (38 cls)  & Avg L$+$D                   & $0{,}278$          & $\mathbf{0{,}556}$ & $\mathbf{0{,}279}$ & $0{,}832$ \\
L2 (38 cls)  & Avg 3                       & $0{,}278$          & $0{,}528$          & $0{,}268$          & $0{,}821$ \\
L2 (38 cls)  & Máx 3                       & $0{,}222$          & $0{,}528$          & $0{,}235$          & $0{,}797$ \\
L2 (38 cls)  & Stack 3                     & $0{,}250$          & $0{,}528$          & $0{,}258$          & $0{,}804$ \\
\hline
L3 (224 cls) & LP PanDerm Base             & $0{,}143$          & ---                & $0{,}132$          & $0{,}735$ \\
L3 (224 cls) & LP PanDerm Large            & $\mathbf{0{,}179}$ & $\mathbf{0{,}286}$ & $\mathbf{0{,}184}$ & $\mathbf{0{,}813}$ \\
L3 (224 cls) & LP DermLIP v2               & $0{,}107$          & $0{,}214$          & $0{,}079$          & $0{,}809$ \\
L3 (224 cls) & Avg L$+$D                   & $\mathbf{0{,}179}$ & $0{,}250$          & $\mathbf{0{,}184}$ & $0{,}802$ \\
L3 (224 cls) & Avg 3                       & $0{,}143$          & $0{,}250$          & $0{,}152$          & $0{,}795$ \\
L3 (224 cls) & Máx 3                       & $0{,}143$          & $0{,}250$          & $0{,}132$          & $0{,}780$ \\
L3 (224 cls) & Stack 3                     & $0{,}143$          & $0{,}250$          & $0{,}158$          & $0{,}791$ \\
\hline
\end{tabular}
\end{table}

El ganador depende del nivel ontológico: ningún codificador lidera
simultáneamente L1, L2 y L3. Sobre L1 lidera la combinación avg
Large$+$DermLIP en AUROC ($0{,}796 \to 0{,}817$). Sobre L2 lidera
DermLIP v2 individual en todas las métricas ($+6{,}7$\,pp de AUROC
sobre PanDerm Large, $0{,}793 \to 0{,}860$), sin que ningún
\emph{ensemble} lo supere; el preentrenamiento contrastivo
texto-imagen invierte la jerarquía justo en el nivel clínico
intermedio, en coherencia con el experimento \emph{zero-shot}
(sección~\ref{sec:dermapixel-zs}). Sobre L3 lidera PanDerm Large
individual, donde la cola larga del nivel limita la señal
complementaria entre codificadores. No existe, por tanto, un
codificador óptimo universal: la mejor opción depende del nivel
ontológico. La lectura de diseño se desarrolla en el
capítulo~\ref{cap:discusion}
(sección~\ref{sec:disc-especializacion}).


\subsection{Ajuste fino parcial del codificador}
\label{sec:dermapixel-ft}
\label{sec:dermapixel-ft-l1}

La sección~\ref{sec:ft-results} cuantifica el efecto del ajuste
fino supervisado (FT) sobre cuatro datasets externos, con
incrementos de hasta $+47{,}1$\,pp BAcc sobre HAM10000. La presente
subsección replica el protocolo sobre DermapixelAI~1.0. El encoder
es PanDerm Large con las dos últimas capas descongeladas
($25{,}2$\,M parámetros entrenables, $8{,}3\%$ del encoder),
seguido de una cabeza \emph{fully-connected}, entrenado con entropía
cruzada reponderada por el desbalance de \emph{Genodermatosis}. La
receta de optimización completa se detalla en el
Anexo~\ref{cap:anexoa} (tabla~\ref{tab:anexoa-ft-recipe}).

Sobre el nivel L1, el ajuste fino parcial alcanza BAcc $0{,}622$
(IC95\% [$0{,}510$--$0{,}735$]) y AUROC $0{,}715$, lo que mejora en
$+18{,}2$\,pp la cifra de validación cruzada agrupada por caso del
sondeo lineal congelado (BAcc $0{,}440 \pm 0{,}058$,
tabla~\ref{tab:dermapixel-cv}), con intervalos no solapados: el
ajuste fino aporta valor real más allá del ruido de partición. No
obstante, no supera a la combinación lineal de codificadores
congelados \texttt{avg\_large\_dermlip}
(sección~\ref{sec:dermapixel-ensemble}; BAcc $0{,}702$, AUROC
$0{,}817$). La replicación sobre L2 y L3 mantiene el mismo
protocolo, variando únicamente la cabeza supervisada. La
tabla~\ref{tab:dermapixel-ft-multilevel} compara las métricas
sobre los tres niveles.

\begin{table}[H]
\centering
\footnotesize
\setlength{\tabcolsep}{4pt}
\renewcommand{\arraystretch}{0.95}
\caption{Ajuste fino parcial del codificador (PanDerm Large, dos
últimas capas descongeladas) sobre los tres niveles ontológicos de
DermapixelAI~1.0. Test fijo $N_{\mathrm{test}} = 36$ en L1/L2 y
$N_{\mathrm{test}} = 28$ en L3. La cabeza L2 de este protocolo
dimensiona $1\,024 \to 37$: son las $37$ clases L2 efectivas en el
\emph{train} del ajuste fino, de las $38$ del corpus tras fusionar
las dos variantes ortográficas de ``queratinización'' ($39$ en
bruto); las $38$ clases son la cardinalidad nominal del nivel L2.
Entre corchetes, IC95\% por \emph{bootstrap}
estratificado ($1\,000$ remuestreos). Mejor valor por columna en
negrita.}
\label{tab:dermapixel-ft-multilevel}
\begin{tabular}{lrrrr}
\hline
Nivel        & Acc@1                            & Acc@3                            & BAcc                             & AUROC \\
\hline
L1 (4 cls)   & $\mathbf{0{,}583}$ [$0{,}444$--$0{,}722$] & $\mathbf{0{,}972}$ [$0{,}917$--$1{,}000$] & $\mathbf{0{,}622}$ [$0{,}510$--$0{,}735$] & $0{,}715$ [$0{,}603$--$0{,}830$] \\
L2 (37 cls)  & $0{,}250$ [$0{,}167$--$0{,}333$] & $0{,}500$ [$0{,}389$--$0{,}611$] & $0{,}324$ [$0{,}265$--$0{,}397$] & $\mathbf{0{,}852}$ [$0{,}814$--$0{,}888$] \\
L3 (224 cls) & $0{,}179$ [$0{,}143$--$0{,}214$] & $0{,}286$ [$0{,}214$--$0{,}357$] & $0{,}184$ [$0{,}158$--$0{,}211$] & $0{,}804$ [$0{,}773$--$0{,}835$] \\
\hline
\end{tabular}
\end{table}

El FT denso parcial mejora al sondeo lineal en L1 y L2, pero no es
la forma de adaptación más eficiente, como revela la comparación
directa con la adaptación de bajo rango.

\subsection{LoRA frente a ajuste fino denso}

Sobre el nivel L2, el FT parcial denso ($25{,}2$\,M parámetros
entrenables) alcanza BAcc $0{,}324$ y AUROC $0{,}852$, cifras que la
adaptación de bajo rango (LoRA~\cite{lora}) del protocolo Dermapixel
R0 (sección~\ref{sec:dermapixel-spanderm-v0}) supera ---BAcc
$0{,}363 \pm 0{,}007$ y AUROC $0{,}855 \pm 0{,}006$--- movilizando
$48$ veces menos parámetros
(tabla~\ref{tab:dermapixel-ft-vs-lora}).

\begin{table}[H]
\centering
\footnotesize
\setlength{\tabcolsep}{4pt}
\renewcommand{\arraystretch}{0.95}
\caption{Comparativa sobre el nivel L2 de DermapixelAI~1.0 entre el
sondeo lineal congelado (referencia), el ajuste fino parcial denso
con dos últimas capas descongeladas y el protocolo Dermapixel R0
(cabeza supervisada L2) con
adaptación LoRA. Mismo conjunto de test ($N_{\mathrm{test}} = 36$,
$37$--$38$ clases según consolidación de variantes ortográficas
de queratinización descrita en
sección~\ref{sec:dermapixel-spanderm-v0}). Mejor valor por
columna en negrita.}
\label{tab:dermapixel-ft-vs-lora}
\begin{tabular}{lrrr}
\hline
Protocolo                       & BAcc L2                         & AUROC L2                       & Params entren. \\
\hline
LP test fijo (referencia, tabla~\ref{tab:dermapixel-lp}) & $0{,}265$                       & $0{,}793$                      & $39\,K$ \\
FT parcial denso (este trabajo) & $0{,}324$                       & $0{,}852$                      & $25{,}2$\,M \\
\textbf{Dermapixel R0 (LoRA)}   & $\mathbf{0{,}363 \pm 0{,}007}$  & $\mathbf{0{,}855 \pm 0{,}006}$ & $\mathbf{524\,K}$ \\
\hline
\end{tabular}
\end{table}

El experimento muestra que, bajo el régimen de datos disponible
($N_{\mathrm{train}} = 874$), el cuello de botella no es la
capacidad del modelo sino la eficiencia de la adaptación: LoRA
obtiene mejores métricas que el ajuste fino denso parcial
movilizando únicamente el $2\%$ de sus parámetros entrenables
($524\,K$ frente a $25{,}2$\,M). Este resultado justifica la
decisión arquitectónica de Dermapixel R0
(sección~\ref{sec:dermapixel-spanderm-v0}).

\subsection{Limitaciones del ajuste fino en la cola larga}

Sobre el nivel L3, el FT denso parcial alcanza BAcc $0{,}184$ y
AUROC $0{,}804$ sobre las $224$ clases efectivas, sin superar al
sondeo lineal congelado sobre el mismo \emph{split} (BAcc $0{,}184$,
AUROC $0{,}813$). La cola larga severa del nivel ($77\%$ de las
clases con $\leq 5$ ejemplos) anula el beneficio del ajuste fino
denso, lo que reserva L3 al ranking por prototipos coseno de
Dermapixel R0 (tabla~\ref{tab:dermapixel-spanderm-l3}) o a la
ampliación del dataset vía banco hospitalario del HUSLL.

Tres \emph{caveats} acotan la interpretación de las cifras de esta
sección: el conjunto de test es reducido ($N = 36$ en L1/L2,
$N = 28$ en L3), lo que ensancha los IC95\%; la categoría etiológica
\emph{Genodermatosis} no aparece en el test de L1 ($0 / 36$), por lo
que las métricas reflejan una clasificación efectiva de tres clases;
y la cola larga de L3 limita el rendimiento absoluto alcanzable bajo
cualquier protocolo de ajuste denso sobre el dataset actual.

\section{Validación \emph{zero-shot} multimodal sobre DermapixelAI 1.0}
\label{sec:dermapixel-zs}

La pregunta no es qué modelo obtiene la mejor accuracy, sino hasta
qué punto un modelo fundacional dermatológico transfiere a un
dataset clínico en castellano sin entrenar ninguna cabeza
supervisada. El hallazgo principal es que la brecha entre el régimen
\emph{zero-shot} y el aprendizaje supervisado depende del nivel
ontológico: permanece amplia en L1 y L3, pero prácticamente
desaparece en L2, donde DermLIP v2 \emph{zero-shot} alcanza una
AUROC ($0{,}794$) indistinguible de la del sondeo lineal supervisado
de PanDerm Large ($0{,}793$).

El experimento opera sobre el mismo conjunto de test que la
sección~\ref{sec:dermapixel-results} ($N_{\mathrm{test}} = 36$ en
L1/L2, $N_{\mathrm{test}} = 28$ en L3), sin entrenar ninguna cabeza
clasificadora y conservando la garantía de no solapamiento con
Derm1M ($0 / 507\,657$ coincidencias por hash MD5). Se evalúan tres
modelos vision-lenguaje contrastivos: DermLIP v2 (preentrenado sobre
Derm1M), SigLIP-SO400M (generalista, WebLI) y BiomedCLIP
(PubMed-15M), contrastando además la formulación del \emph{prompt}
en castellano frente a inglés. La predicción se obtiene por
\emph{argmax} sobre la similitud coseno con un \emph{ensemble} de
tres plantillas textuales por nivel; el nivel L3 se evalúa
únicamente en castellano, al no disponer el vocabulario ontológico
de traducción inglesa validada para sus $224$ clases. La
tabla~\ref{tab:dermapixel-zs} reporta los resultados.

\begin{table}[H]
\centering
\footnotesize
\setlength{\tabcolsep}{4pt}
\renewcommand{\arraystretch}{0.95}
\caption{Clasificación \emph{zero-shot} multimodal sobre el
conjunto de test de DermapixelAI~1.0 ($N_{\mathrm{test}} = 36$ en
L1/L2, $N_{\mathrm{test}} = 28$ en L3). Entre corchetes, IC95\%
por \emph{bootstrap} estratificado ($1\,000$ remuestreos) para
las dos métricas primarias (Acc@1 y AUROC). Las métricas BAcc y
Acc@3 se reportan como valor puntual por concisión; los IC95\%
correspondientes están disponibles en el material reproducible.
L3 se evalúa únicamente en castellano. Mejor valor por columna,
nivel e idioma en negrita.}
\label{tab:dermapixel-zs}
\begin{tabular}{lllrrrr}
\hline
Nivel & Modelo & Lang & Acc@1 & Acc@3 & BAcc & AUROC \\
\hline
L1 (4 cls)   & BiomedCLIP    & es & $\mathbf{0{,}389}$ [$0{,}250$--$0{,}528$]
                                  & $\mathbf{0{,}944}$
                                  & $0{,}375$
                                  & $0{,}679$ [$0{,}562$--$0{,}781$] \\
L1 (4 cls)   & BiomedCLIP    & en & $0{,}389$ [$0{,}250$--$0{,}528$]
                                  & $\mathbf{0{,}972}$
                                  & $0{,}324$
                                  & $0{,}639$ [$0{,}524$--$0{,}738$] \\
L1 (4 cls)   & SigLIP-SO400M & es & $0{,}222$ [$0{,}111$--$0{,}361$]
                                  & $0{,}917$
                                  & $0{,}178$
                                  & $0{,}650$ [$0{,}571$--$0{,}735$] \\
L1 (4 cls)   & SigLIP-SO400M & en & $0{,}333$ [$0{,}194$--$0{,}472$]
                                  & $0{,}889$
                                  & $0{,}292$
                                  & $0{,}663$ [$0{,}595$--$0{,}719$] \\
L1 (4 cls)   & DermLIP v2    & es & $0{,}361$ [$0{,}222$--$0{,}500$]
                                  & $0{,}861$
                                  & $0{,}379$
                                  & $0{,}699$ [$0{,}600$--$0{,}794$] \\
L1 (4 cls)   & DermLIP v2    & en & $\mathbf{0{,}500}$ [$0{,}333$--$0{,}667$]
                                  & $0{,}917$
                                  & $\mathbf{0{,}441}$
                                  & $\mathbf{0{,}709}$ [$0{,}565$--$0{,}841$] \\
\hline
L2 (38 cls)  & BiomedCLIP    & es & $0{,}111$ [$0{,}111$--$0{,}111$]
                                  & $0{,}278$
                                  & $0{,}118$
                                  & $0{,}647$ [$0{,}604$--$0{,}688$] \\
L2 (38 cls)  & BiomedCLIP    & en & $0{,}083$ [$0{,}056$--$0{,}111$]
                                  & $0{,}250$
                                  & $0{,}088$
                                  & $0{,}655$ [$0{,}601$--$0{,}713$] \\
L2 (38 cls)  & SigLIP-SO400M & es & $0{,}139$ [$0{,}139$--$0{,}139$]
                                  & $0{,}250$
                                  & $0{,}235$
                                  & $0{,}636$ [$0{,}589$--$0{,}680$] \\
L2 (38 cls)  & SigLIP-SO400M & en & $0{,}194$ [$0{,}111$--$0{,}278$]
                                  & $0{,}278$
                                  & $0{,}221$
                                  & $0{,}697$ [$0{,}640$--$0{,}747$] \\
L2 (38 cls)  & DermLIP v2    & es & $\mathbf{0{,}222}$ [$0{,}167$--$0{,}278$]
                                  & $0{,}333$
                                  & $0{,}265$
                                  & $0{,}762$ [$0{,}721$--$0{,}805$] \\
L2 (38 cls)  & DermLIP v2    & en & $\mathbf{0{,}250}$ [$0{,}194$--$0{,}306$]
                                  & $\mathbf{0{,}417}$
                                  & $\mathbf{0{,}294}$
                                  & $\mathbf{0{,}794}$ [$0{,}747$--$0{,}837$] \\
\hline
L3 (224 cls) & BiomedCLIP    & es & $0{,}107$ [$0{,}036$--$0{,}179$]
                                  & $0{,}214$
                                  & $0{,}079$
                                  & $0{,}703$ [$0{,}673$--$0{,}732$] \\
L3 (224 cls) & SigLIP-SO400M & es & $0{,}071$ [$0{,}000$--$0{,}143$]
                                  & $0{,}214$
                                  & $0{,}053$
                                  & $0{,}696$ [$0{,}667$--$0{,}723$] \\
L3 (224 cls) & DermLIP v2    & es & $\mathbf{0{,}107}$ [$0{,}071$--$0{,}143$]
                                  & $\mathbf{0{,}214}$
                                  & $\mathbf{0{,}079}$
                                  & $\mathbf{0{,}740}$ [$0{,}719$--$0{,}763$] \\
\hline
\end{tabular}
\end{table}

Tres hallazgos cuantitativos estructuran la interpretación.

\textit{Primero}, DermLIP v2 obtiene el mejor o equivalente
rendimiento en las cinco columnas métricas sobre los tres niveles
ontológicos, con excepción de Acc@3 L1 (BiomedCLIP en castellano
e inglés alcanza $0{,}944$ y $0{,}972$ respectivamente, frente a
$0{,}861$ y $0{,}917$ de DermLIP v2). La superioridad
sistemática es coherente con su preentrenamiento sobre Derm1M, el
único dataset de los tres específicamente dermatológico (un millón
de pares imagen-texto curados sobre la familia
PanDerm)~\cite{derm1m}.

\textit{Segundo}, la formulación del \emph{prompt} en castellano
penaliza el rendimiento de los modelos preentrenados sobre dataset
mayoritariamente anglófonos. Sobre L1 Acc@1, DermLIP v2 cae de
$0{,}500$ (inglés) a $0{,}361$ (castellano), una pérdida de
$13{,}9$\,pp; SigLIP-SO400M cae de $0{,}333$ a $0{,}222$
($-11{,}1$\,pp). BiomedCLIP, en cambio, mantiene
$0{,}389$ en ambos idiomas (\emph{gap} nulo). La penalización
observada es consistente con el predominio del inglés en los corpus
textuales utilizados durante el preentrenamiento, aunque el diseño
experimental no permite atribuir causalmente la diferencia a este
único factor. La penalización cuantificada motiva la
construcción de modelos vision-lenguaje contrastivos con rama
textual en castellano, línea futura abierta por el dataset
DermapixelAI~1.0.

\textit{Tercero}, la comparación con el sondeo lineal sobre el
mismo conjunto de test (tabla~\ref{tab:dermapixel-lp}) cuantifica
la brecha entre un sistema entrenado y un sistema
\emph{zero-shot}. Sobre L1, PanDerm Large bajo LP alcanza
BAcc $= 0{,}735$ frente a $0{,}441$ del mejor \emph{zero-shot}
(DermLIP v2 en inglés), una diferencia de $+29{,}4$\,pp a favor
del sondeo lineal. Sobre L3, la brecha es de $+10{,}5$\,pp BAcc
($0{,}184$ vs $0{,}079$) y $+7{,}3$\,pp AUROC ($0{,}813$ vs
$0{,}740$). El comportamiento se invierte parcialmente sobre L2:
LP PanDerm Large obtiene AUROC $0{,}793$ y \emph{zero-shot}
DermLIP v2 en inglés alcanza $0{,}794$, un empate efectivo dentro
del margen de los intervalos de confianza. La interpretación es
relevante: las \emph{features} pre-entrenadas de la familia
PanDerm/DermLIP son ya competitivas a nivel de subcategoría
clínica L2 sin necesidad de cabeza entrenada, mientras que la
discriminación etiológica de grano grueso (L1) y de cola larga
diagnóstica (L3) sí requiere un clasificador supervisado.
Adicionalmente, Acc@3 sobre L1 oscila entre $0{,}861$ (DermLIP v2
en castellano) y $0{,}972$ (BiomedCLIP en inglés), magnitudes
operativamente relevantes para escenarios de cribado en los que
el sistema presente las tres categorías etiológicas más probables
al clínico antes del juicio diagnóstico.

La tabla~\ref{tab:dermapixel-zs-gaps} resume las brechas
cuantificadas en esta sección.

\begin{table}[H]
\centering
\footnotesize
\setlength{\tabcolsep}{4pt}
\renewcommand{\arraystretch}{0.95}
\caption{Brechas \emph{zero-shot} cuantificadas sobre DermapixelAI~1.0:
penalización por idioma (castellano frente a inglés) y distancia al
sondeo lineal supervisado (LP).}
\label{tab:dermapixel-zs-gaps}
\begin{tabular}{llr}
\hline
Contraste                                & Métrica & \emph{Gap} \\
\hline
DermLIP v2, L1, castellano vs.\ inglés   & Acc@1 & $-13{,}9$\,pp \\
SigLIP-SO400M, L1, castellano vs.\ inglés & Acc@1 & $-11{,}1$\,pp \\
DermLIP v2, L2, castellano vs.\ inglés   & AUROC & $-3{,}2$\,pp \\
LP supervisado vs.\ \emph{zero-shot}, L1 & BAcc  & $+29{,}4$\,pp \\
LP supervisado vs.\ \emph{zero-shot}, L3 & BAcc  & $+10{,}5$\,pp \\
\hline
\end{tabular}
\end{table}

\subsection*{Caveats}

Las cautelas sobre el tamaño del conjunto de test ($N = 36$ en
L1/L2, $N = 28$ en L3) y sobre la ausencia de solapamiento frente
a Derm1M ya enunciadas para el sondeo lineal
(sección~\ref{sec:dermapixel-results}) aplican también a esta
evaluación zero-shot. Dos \emph{caveats} adicionales son
específicos del régimen zero-shot. Primero, el \emph{ensemble} de
tres plantillas textuales por nivel constituye una formulación
mínima de \emph{prompt engineering}; la literatura sobre
\emph{zero-shot} multimodal emplea típicamente entre diez y ochenta
plantillas generadas con asistencia de LLMs, lo que sitúa las
cifras reportadas como cota inferior del rendimiento alcanzable.
Segundo, la evaluación L3 únicamente en castellano impide
cuantificar la penalización por idioma sobre la cola larga
diagnóstica, dado que el vocabulario ontológico no dispone de
traducción inglesa validada para sus $224$ clases.

\section{Ajuste fino con LoRA: Dermapixel R0 (cabeza supervisada L2)}
\label{sec:dermapixel-spanderm-v0}

Los experimentos anteriores evalúan modelos existentes sobre
DermapixelAI~1.0 como extractores de \emph{features} congelados con
cabezas ligeras o clasificación contrastiva \emph{zero-shot}. La
presente sección aborda la contribución propia: adaptar PanDerm
Large al contexto clínico en castellano mediante LoRA. Dermapixel
R0, con menos del $0{,}2\%$ de parámetros entrenables, alcanza la
mejor discriminación de subcategorías documentada sobre el dataset;
sus métricas L2 y el ranking L3 se reportan a continuación. La
síntesis de las técnicas de \emph{transfer learning} ejecutadas y
pendientes sobre el dataset se recoge en la
sección~\ref{sec:dermapixel-futuro} del
capítulo~\ref{cap:conclusiones}.

La sección~\ref{sec:dermapixel-ft-l1} documenta el ajuste fino
parcial denso sobre L1, que mejora la BAcc en $+18{,}2$\,pp sobre
la cifra de validación cruzada del sondeo lineal congelado pero
no supera al \emph{ensemble} multi-codificador
(sección~\ref{sec:dermapixel-ensemble}). La presente
sección introduce \textbf{Dermapixel R0}, el primer modelo
fundacional dermatológico en español propuesto en este trabajo,
inicializado desde el codificador PanDerm Large y adaptado sobre el
dataset DermapixelAI~1.0 mediante adaptación de bajo rango
(\emph{Low-Rank Adaptation}, LoRA)~\cite{lora}. La denominación R0
designa una versión base, un punto de partida entrenado sobre un
corpus reducido, sobre el que sentar las bases de iteraciones
posteriores. La presente
sección reporta la cabeza supervisada L2 de Dermapixel R0 sobre el
nivel L2 ($38$ clases efectivas tras la consolidación de
variantes ortográficas de queratinización ya descrita en la
sección~\ref{sec:dermapixel-ft-l1}) y el ranking L3 por
prototipos sobre las $250$ clases efectivas del dataset. La
elección de L2 como objetivo de la cabeza supervisada responde a
la cardinalidad operativa del nivel (rendimiento clínicamente
informativo sin alcanzar el régimen de cola larga extrema de L3,
con $3{,}5$ muestras/clase en $N_{\mathrm{train}} = 874$). El
ranking L3 sigue la estrategia prevista en el diseño v0: prototipos
calculados como la media de las representaciones del conjunto de
entrenamiento por clase L3 y predicción por similitud coseno. La
configuración completa del protocolo (encoder congelado, adaptación
LoRA $r=16$, optimización, augmentaciones y semillas) se detalla en
el Anexo~\ref{cap:anexoa}
(tabla~\ref{tab:dermapixel-spanderm-config}).

Las cifras reportadas a continuación corresponden a la media y
la desviación estándar sobre las tres semillas
($\{42, 43, 44\}$). La tabla~\ref{tab:dermapixel-spanderm-l2}
reporta las dos estimaciones complementarias sobre el conjunto de
test L2 ($N_{\mathrm{test}} = 36$): la cifra del \emph{best val
checkpoint} y la media de las últimas cinco épocas sin selección
de \emph{checkpoint}. La segunda estimación caracteriza el
comportamiento asintótico del entrenamiento sin la varianza
asociada al criterio de selección.

\begin{table}[H]
\centering
\footnotesize
\setlength{\tabcolsep}{4pt}
\renewcommand{\arraystretch}{0.95}
\caption{Métricas de Dermapixel R0 (cabeza supervisada L2) sobre el
conjunto de test L2 de
DermapixelAI~1.0 ($N_{\mathrm{test}} = 36$, $38$ clases con
queratinización consolidada). Cifras: media $\pm$ desviación
estándar sobre tres semillas ($\{42, 43, 44\}$). La columna
\emph{best val} aplica selección por máxima BAcc de validación;
la columna \emph{últimas $5$ ép.} promedia las cinco épocas
finales sin selección. Mejor por fila en negrita.}
\label{tab:dermapixel-spanderm-l2}
\begin{tabular}{lrr}
\hline
Métrica   & \emph{Best val \emph{checkpoint}} & Media últimas $5$ ép. \\
\hline
Acc@1     & $\mathbf{0{,}250 \pm 0{,}000}$ & $0{,}209 \pm 0{,}005$ \\
Acc@3     & $\mathbf{0{,}500 \pm 0{,}023}$ & $0{,}457 \pm 0{,}011$ \\
BAcc      & $\mathbf{0{,}363 \pm 0{,}007}$ & $0{,}309 \pm 0{,}006$ \\
AUROC     & $\mathbf{0{,}855 \pm 0{,}006}$ & $0{,}837 \pm 0{,}008$ \\
W-F1      & $\mathbf{0{,}218 \pm 0{,}013}$ & --- \\
Kappa     & $\mathbf{0{,}220 \pm 0{,}001}$ & --- \\
\hline
\end{tabular}
\end{table}

La estabilidad entre semillas es alta sobre todas las métricas
reportadas (desviación estándar $\leq 0{,}013$ en valor absoluto
y $\leq 0{,}008$ sobre AUROC y BAcc). La diferencia entre el
\emph{best val checkpoint} y la media de las últimas cinco épocas
($+5{,}4$\,pp BAcc, $+1{,}8$\,pp AUROC, $+4{,}1$\,pp Acc@1)
indica que la curva de entrenamiento exhibe un patrón de
sobreajuste suave análogo al observado en el FT parcial denso de
la sección~\ref{sec:dermapixel-ft-l1}, pero atenuado por la
restricción al $0{,}17\%$ de parámetros entrenables que impone
la adaptación LoRA.

La tabla~\ref{tab:dermapixel-spanderm-comp} sitúa las cifras de
Dermapixel R0 frente a las dos referencias internas sobre el mismo
conjunto de test y el mismo nivel ontológico: el sondeo lineal
congelado bajo validación cruzada agrupada por caso
(tabla~\ref{tab:dermapixel-cv}), el sondeo lineal sobre el split
fijo (tabla~\ref{tab:dermapixel-lp}) y el codificador DermLIP v2
individual reportado en la sección~\ref{sec:dermapixel-ensemble}
como mejor opción sobre L2 sin entrenamiento de cabeza.

\begin{table}[H]
\centering
\footnotesize
\setlength{\tabcolsep}{4pt}
\renewcommand{\arraystretch}{0.95}
\caption{Comparativa sobre el nivel L2 de DermapixelAI~1.0
($38$ clases con queratinización consolidada,
$N_{\mathrm{test}} = 36$) entre Dermapixel R0 y las referencias
internas del trabajo. La fila LP CV se considera la cifra honesta
del sondeo lineal congelado por su independencia del \emph{split}
fijo (sección~\ref{sec:dermapixel-results}). Mejor valor por
columna en negrita.}
\label{tab:dermapixel-spanderm-comp}
\adjustbox{max width=\textwidth}{%
\begin{tabular}{lllrr}
\hline
Método                                              & Encoder                       & Params entren.    & BAcc L2                         & AUROC L2 \\
\hline
LP test fijo (tabla~\ref{tab:dermapixel-lp})        & PanDerm Large $1\,024$ (cong.) & $39\,K$           & $0{,}265$                       & $0{,}793$ \\
LP $5$-fold CV (tabla~\ref{tab:dermapixel-cv})      & PanDerm Large $1\,024$ (cong.) & $39\,K$           & $0{,}168 \pm 0{,}020$           & $0{,}749 \pm 0{,}019$ \\
DermLIP v2 individual (sección~\ref{sec:dermapixel-ensemble}) & DermLIP $512$ (cong.)        & $19{,}5\,K$       & $0{,}279$                       & $\mathbf{0{,}860}$ \\
\textbf{Dermapixel R0} (este trabajo)               & PanDerm Large $1\,024$ $+$ LoRA & $524\,K$          & $\mathbf{0{,}363 \pm 0{,}007}$  & $0{,}855 \pm 0{,}006$ \\
\hline
\end{tabular}}
\end{table}

La comparativa documenta dos hallazgos sobre el mismo conjunto de
test. Dermapixel R0 mejora la BAcc L2 en $+19{,}5$\,pp respecto a
la validación cruzada del sondeo lineal congelado
($0{,}168 \to 0{,}363$) y en $+8{,}4$\,pp respecto a DermLIP v2
individual ($0{,}279 \to 0{,}363$). Su AUROC L2
($0{,}855 \pm 0{,}006$) empata estadísticamente con la de DermLIP
v2 ($0{,}860$). El compromiso se interpreta como
complementariedad: la cabeza supervisada con LoRA sobre PanDerm
Large discrimina mejor las clases minoritarias (mejor BAcc),
mientras que el preentrenamiento contrastivo retiene una ligera
ventaja sobre la AUROC \emph{one-vs-rest} global.

El nivel L3 se aborda mediante un ranking por prototipos coseno
---la media de las representaciones de entrenamiento por clase, con
predicción por similitud coseno---, estrategia prevista para mitigar
la cola larga del nivel sin una cabeza supervisada directa, inviable
bajo $3{,}5$ muestras/clase de media. El ranking alcanza Acc@1
$= 0{,}167$ y Acc@5 $= 0{,}296$, equivalentes a $\times 42$ y
$\times 15$ sobre el azar respectivamente, pero permanece lejos del
régimen operativo clínico: la cola larga del vocabulario ($178$ de
$250$ clases con $\leq 5$ imágenes en \emph{train}) produce
prototipos ruidosos. La tabla completa figura en el
Anexo~\ref{cap:anexoj} (tabla~\ref{tab:dermapixel-spanderm-l3}).

Dos hallazgos cierran la sección. \textit{Primero}, Dermapixel R0
con sólo $0{,}17\%$ de parámetros entrenables iguala o supera a las
tres referencias internas sobre BAcc L2 ($0{,}363 \pm 0{,}007$
frente a $0{,}168 \pm 0{,}020$ del LP CV, $0{,}265$ del LP fijo y
$0{,}279$ de DermLIP v2 individual), la mejor BAcc L2 documentada
sobre el dataset, lo que valida la adaptación de bajo rango como
protocolo preferente en castellano. \textit{Segundo}, su AUROC L2
($0{,}855 \pm 0{,}006$) empata estadísticamente con la de DermLIP
v2 nativo ($0{,}860$): la adaptación textual ligera de PanDerm
Large mediante LoRA alcanza el mismo techo discriminativo que el
preentrenamiento contrastivo sobre Derm1M. La progresión hacia una
cabeza L3 directa queda condicionada a la incorporación del banco
hospitalario del HUSLL (umbral $\geq 30$ imágenes/clase L3,
sección~\ref{sec:dermapixel-futuro}).

\section{Dermapixel R0: rama contrastiva multimodal castellano
zero-shot sobre DermapixelAI~1.0}
\label{sec:spanderm-clip}

A diferencia de la cabeza supervisada L2, que fija las clases por
anticipado, esta rama mapea texto clínico libre en castellano sobre
el espacio de imágenes sin vocabulario cerrado, condición relevante
para la búsqueda documental sobre el archivo clínico. La progresión
de cinco iteraciones (v1$\to$v5) revela que la palanca decisiva es
la capacidad del codificador visual, no la escala de datos externos,
e incluye un resultado negativo de escala.

La presente sección cierra una configuración no cubierta por las
anteriores: un sistema vision-lenguaje contrastivo \emph{dual-encoder}
con encoder textual multilingüe entrenado sobre el propio dataset
DermapixelAI~1.0, evaluado en \emph{zero-shot} \emph{open-class}
sobre el vocabulario completo de la ontología (L1, L2 y L3) en
castellano. Frente a DermLIP, sustituye el encoder textual
anglosajón (PubMedBERT) por un encoder multilingüe adaptado al
lenguaje clínico castellano del dataset. La justificación
cuantitativa frente a un \emph{wrapper} con modelo fundacional
inglés más traducción ES--EN se desarrolla en el
Anexo~\ref{cap:anexoi}.

\subsection{Arquitectura, datos y protocolo experimental}

La rama contrastiva de Dermapixel R0 sigue el patrón dual-encoder
estilo CLIP~\cite{clip2021}: un encoder visual de
$1\,024$ dimensiones (PanDerm Large, ViT-L/16, CLS \emph{pooling}),
un encoder textual multilingüe de $768$ dimensiones
(\texttt{intfloat/multilingual-e5-base}), capas de proyección
lineales con normalización L2 a un espacio compartido de
$512$ dimensiones y entrenamiento por pérdida InfoNCE bidireccional
(imagen$\to$texto y texto$\to$imagen) con temperatura aprendible.
El dataset de entrenamiento es DermapixelAI~1.0
(Anexo~\ref{cap:anexoc}) bajo el \emph{split} per-case estratificado
por L1 ($854$ \emph{train} / $107$ \emph{val} / $101$ \emph{test}
imágenes; $521 / 66 / 66$ casos), idéntico al utilizado en las
secciones~\ref{sec:dermapixel-results} y~\ref{sec:dermapixel-zs}.
Los conjuntos de validación y test quedan fijados por un manifiesto
SHA-256 en cada \emph{checkpoint} y no se tocan en ninguna iteración
del experimento. Para amplificar la señal textual del dataset se
generan, mediante un modelo de lenguaje grande a partir del texto de
cada caso y sus metadatos, ocho perspectivas clínicas por imagen
(morfología, dermatoscopia, criterios ABCDE, diagnóstico
diferencial, informe clínico, etc.) que, combinadas con nueve
variantes textuales base, producen del orden de $22\,000$ pares
(imagen, texto). Cada época muestra cada imagen con una variante
textual seleccionada por hash determinista, lo que garantiza
reproducibilidad bajo reanudación. Los detalles de generación
(coste, longitud media y procedencia) figuran en el
Anexo~\ref{cap:anexoi}.

Se entrenan cinco iteraciones bajo configuraciones crecientes de
capacidad arquitectónica y, en la última, bajo ampliación del
dataset. La progresión aísla por construcción tres palancas
independientes: adaptación del lado textual (v1$\to$v2, LoRA
textual sobre el encoder E5 multilingüe), adaptación del lado
visual (v2$\to$v3 con LoRA $r=16$ sobre los últimos cuatro bloques
del ViT, y v3$\to$v4 ampliando a $r=32$ sobre los últimos doce) y
escala del dataset (v4$\to$v5, $15\,128$ imágenes externas
genéricas con \emph{captions} sintéticas bajo idéntica
configuración LoRA que v4). La configuración completa de las cinco
iteraciones se detalla en el Anexo~\ref{cap:anexoa}
(tabla~\ref{tab:spanderm-config}); los tiempos de entrenamiento y
la lección metodológica sobre la longitud de las \emph{captions}
sintéticas figuran en el Anexo~\ref{cap:anexoi}. El objetivo
experimental, fijado a priori, es L3 \emph{accuracy}@1
$\geq 0{,}25$ en \emph{zero-shot} sobre las $54$ clases L3
representadas en el conjunto de test.

\subsection{Resultados \emph{zero-shot}: progresión metodológica
v1$\to$v4 y ablación de escala v5}

La tabla~\ref{tab:spanderm-zs} reporta la \emph{accuracy}@1 sobre
las $101$ imágenes del conjunto de test bajo régimen híbrido
\emph{H6} (combinación lineal de logits visuales y textuales con
pesos $w_v = 0{,}3$ y $w_t = 0{,}7$, calibrados sobre validación),
aplicado de forma uniforme a las cinco iteraciones para garantizar
la comparativa cross-checkpoint directa. Los niveles L1 y L2
\emph{zero-shot} sólo se reportan para v4 y v5, los
\emph{checkpoints} relevantes para los regímenes operativos del
prototipo.

\begin{table}[H]
\centering
\footnotesize
\setlength{\tabcolsep}{3pt}
\renewcommand{\arraystretch}{0.95}
\caption{Clasificación \emph{zero-shot} de la rama contrastiva de
Dermapixel R0 sobre el
conjunto de test de $101$ imágenes (per-case estratificado por
L1) bajo régimen híbrido \emph{H6} homogéneo
($w_v = 0{,}3$, $w_t = 0{,}7$) aplicado a las cinco iteraciones.
La columna \emph{Objetivo} corresponde al umbral fijado en el
diseño experimental antes de observar resultados. Mejor cifra
por fila excluyendo la columna \emph{Objetivo} en negrita.
Provenance: commit \texttt{d0190c3} del repositorio
\texttt{spanderm-clip}.}
\label{tab:spanderm-zs}
\begin{tabular}{lrrrrrrr}
\hline
Tarea                        & Random    & v1                  & v2                  & v3                  & v4                  & v5                  & Objetivo \\
\hline
L1 acc@1 ($4$ cls)            & $0{,}250$ & ---                 & ---                 & ---                 & $0{,}4257$          & $\mathbf{0{,}4554}$ & ---           \\
L2 acc@1 ($23$ cls)           & $0{,}043$ & ---                 & ---                 & ---                 & $0{,}1837$          & $\mathbf{0{,}2041}$ & ---           \\
\textbf{L3 acc@1 ($54$ cls)}  & $0{,}019$ & $0{,}1782$          & $0{,}1980$          & $0{,}1980$          & $\mathbf{0{,}2475}$ & $0{,}2079$          & $\geq 0{,}25$ \\
\hline
\end{tabular}
\end{table}

La progresión v1$\to$v4 responde al diseño experimental sobre el
nivel L3. Las tres primeras iteraciones quedan estadísticamente
empatadas en el intervalo $[0{,}1782,\,0{,}1980]$ ($\pm 2$
imágenes sobre $N = 101$); la figura~\ref{fig:spanderm-epoch-sweep}
muestra que v3 oscila sin aproximarse al objetivo. La cuarta
iteración alcanza $0{,}2475$, a $0{,}0025$ del umbral objetivo
($0{,}25$): la distancia equivale a una imagen de test, dentro
del ruido de muestreo. El salto $+4{,}95$\,pp entre v3 y v4 bajo
idéntico régimen \emph{H6} se atribuye a la ampliación de la
capacidad arquitectónica del image encoder (LoRA $r=32$ sobre doce
bloques frente a $r=16$ sobre cuatro). El cuello de botella no
estaba, por tanto, en la formulación textual ni en las
\emph{captions} ricas, sino en la capacidad del lado visual para
reesculpir las \emph{features} \emph{fine-grained} castellanas; el
\emph{lift} de recuperación $14{,}0\times$ ya alcanzado por v2 con
LoRA textual ligera (reportado más adelante) confirma que la
alineación cross-modal se logró sin trasladar el \emph{zero-shot}
L3 por encima del techo del encoder visual congelado.

La quinta iteración es una ablación deliberada de la hipótesis de
escala. Bajo idéntica configuración LoRA que v4, la ampliación del
dataset con $15\,128$ imágenes externas genéricas reduce el L3
\emph{accuracy}@1 a $0{,}2079$ ($-4{,}0$\,pp respecto a v4),
resultado \emph{flat-to-negative} discutido mecanísticamente en
la sección~\ref{sec:spanderm-escala-especificidad}: el encoder
visual mejora aisladamente bajo régimen visual puro (\emph{H1}
$0{,}168 \to 0{,}238$) y el techo \emph{Linear Probing} L2 sube
$+54{,}4\%$ relativo, pero el régimen híbrido L3 \emph{fine-grained}
se degrada porque las \emph{captions} externas erosionan la
alineación del registro lingüístico original. La cola larga no se
rellena: las $327$ clases L3 con $\leq 5$ imágenes permanecen sin
cobertura adicional.

\begin{figure}[H]
\centering
\includegraphics[width=0.88\textwidth]{fig_spanderm_epoch_sweep.png}
\caption{Barrido de épocas de la iteración contrastiva v3 de
Dermapixel R0 sobre el
conjunto de test de $101$ imágenes. \emph{Accuracy}@1 L1, L2 y L3
y \texttt{val i2t\_R@5} frente al paso de entrenamiento. L3
oscila en el intervalo $[0{,}04,\,0{,}12]$ a lo largo de las diez
épocas sin aproximarse al objetivo $\geq 0{,}25$, mientras
que \texttt{val i2t\_R@5} alcanza su máximo en la época $2$ y
desciende con el entrenamiento conjunto adicional del codificador
visual.}
\label{fig:spanderm-epoch-sweep}
\end{figure}

\subsection{Recuperación cross-modal: resultado favorable}

El segundo eje de evaluación es la recuperación cross-modal
per-imagen sobre los $107$ pares (imagen, caption rica) del
conjunto de validación, con script idéntico para las tres
versiones. La tabla~\ref{tab:spanderm-retrieval} reporta el
régimen \emph{rich-caption} (campo \texttt{llm\_case\_summary},
el más representativo del lenguaje clínico real de la
Dra.~R.~Taberner).

\begin{table}[H]
\centering
\footnotesize
\setlength{\tabcolsep}{4pt}
\renewcommand{\arraystretch}{0.95}
\caption{Recuperación cross-modal i2t R@$5$ de la rama contrastiva
de Dermapixel R0 sobre
los $107$ pares (imagen, caption rica) de validación, régimen
\emph{rich-caption} (\texttt{llm\_case\_summary}). La columna
\emph{Lift} reporta el factor multiplicativo sobre el baseline
aleatorio (\emph{pool} $= 107$, baseline $R@5 = 5/107 = 0{,}0467$).
v2 marca el pico histórico de alineación cross-modal castellano;
v5 muestra descenso coherente con la degradación documentada en
la tabla~\ref{tab:spanderm-zs} sobre el régimen híbrido. v1 y v3
no se reportan por insuficiencia de datos en el régimen
homogéneo cross-checkpoint. El \emph{lift} (factor de mejora sobre
el azar) se calcula como
$R@5_{\mathrm{abs}} / (k / N_{\mathrm{pool}})$, con
$k = 5$ y $N_{\mathrm{pool}} = 107$; normaliza respecto a la
dificultad intrínseca del \emph{pool} y constituye la métrica
honesta para comparar entre \emph{checkpoints}.}
\label{tab:spanderm-retrieval}
\begin{tabular}{lrrr}
\hline
Versión & i2t R@$5$ absoluto  & Referencia azar & \emph{Lift} sobre azar   \\
\hline
v2      & $\mathbf{0{,}654}$  & $0{,}047$       & $\mathbf{14{,}0\times}$  \\
v4      & $0{,}579$           & $0{,}047$       & $12{,}4\times$           \\
v5      & $0{,}505$           & $0{,}047$       & $10{,}8\times$           \\
\hline
\end{tabular}
\end{table}

La alineación cross-modal castellano-imagen alcanza su pico
histórico en v2: i2t R@$5 = 0{,}654$ frente a la referencia
aleatoria $0{,}047$ sobre \emph{pool} de $107$, equivalente a un
\emph{lift} (factor de mejora sobre el azar) de
$\mathbf{14{,}0\times}$. Esta cifra constituye
la contribución empírica favorable del experimento bajo régimen
\emph{rich-caption} (figura~\ref{fig:spanderm-retrieval}). Las
iteraciones posteriores conservan \emph{lift}
sustancial (v4 $12{,}4\times$, v5 $10{,}8\times$) pero descienden
respecto al pico, en coherencia con la degradación documentada en
la tabla~\ref{tab:spanderm-zs} sobre régimen híbrido: el
entrenamiento conjunto con LoRA visual añadida (v4) y la posterior
incorporación de \emph{captions} externas (v5) desplazan el
espacio embebido respecto al óptimo \emph{rich-caption} castellano
puro alcanzado por v2. El descenso v2$\to$v5 del \emph{lift}
($14{,}0\times \to 10{,}8\times$) acompaña cuantitativamente la
hipótesis de erosión de la \emph{fingerprint} clínica castellana
discutida en
sección~\ref{sec:spanderm-escala-especificidad}.

\begin{figure}[H]
\centering
\includegraphics[width=0.88\textwidth]{fig_spanderm_retrieval.png}
\caption{Recuperación cross-modal i2t R@$5$ de la rama contrastiva
de Dermapixel R0
sobre $107$ pares de validación, régimen \emph{rich-caption}.
Pico v2 de $\mathbf{14{,}0\times}$ de \emph{lift} sobre el azar
($0{,}047 \to 0{,}654$) atribuible al conjunto LoRA textual
$+$ \emph{captions} sintéticas ricas generadas mediante un modelo de
lenguaje grande. v4 y v5 conservan \emph{lift} sustancial
($12{,}4\times$ y $10{,}8\times$ respectivamente) pero descienden
respecto al pico, en coherencia con la degradación documentada
sobre el régimen híbrido.}
\label{fig:spanderm-retrieval}
\end{figure}

\subsection{Techo de \emph{Linear Probing} y diagnóstico
mecanístico}

La pregunta clave del experimento es por qué ningún cambio en el
lado textual mueve la aguja \emph{zero-shot} L3. La respuesta
surge de la medición rigurosa del techo \emph{Linear Probing}
sobre las \emph{features} $1\,024$-D del encoder visual, ejecutada
con script idéntico para v2 (encoder visual frozen) y v3 (LP en
época $2$ y época $6$). La
tabla~\ref{tab:spanderm-lp-ceiling} reúne los tres puntos de la
comparación.

\begin{table}[H]
\centering
\footnotesize
\setlength{\tabcolsep}{4pt}
\renewcommand{\arraystretch}{0.95}
\caption{Techo \emph{Linear Probing} supervisado sobre
\emph{features} $1\,024$-D del image encoder. Comparativa
\emph{frozen} (v2 e$7$) frente a image-LoRA tras entrenamiento
adicional (v3 época $6$). Métrica: \emph{accuracy}@1. Las
cifras de v3 corresponden a la evaluación intermedia conservada
del Sprint97; la reejecución posterior arroja desviaciones
$< 0{,}01$ absolutas, consistentes con ruido de muestreo
($N = 101$) y sin alterar la lectura ($+0{,}01 \approx$ ruido).}
\label{tab:spanderm-lp-ceiling}
\begin{tabular}{lrrrr}
\hline
Nivel              & LP v2-e$7$ frozen & LP v3 best.pt (e$2$) & LP v3 e$6$ ($+$image train) & $\Delta$ frozen$\to$trained \\
\hline
L1 ($4$ vías)       & $0{,}584$         & $0{,}584$            & $0{,}594$                   & $+0{,}01$                   \\
L2 ($38$ vías)      & $0{,}214$         & $0{,}214$            & $0{,}225$                   & $+0{,}01$                   \\
L3 ($225$ vías)     & $0{,}079$         & $0{,}079$            & $0{,}089$                   & $+0{,}01$                   \\
\hline
\end{tabular}
\end{table}

Dos observaciones estructuran el diagnóstico. En
primer lugar, el \emph{zero-shot} L3 alcanzado por v1/v2/v3
($\sim 0{,}13$) supera ya el techo LP supervisado
($0{,}079$--$0{,}089$) sobre el mismo encoder visual frozen.
Las \emph{features} $1\,024$-D del PanDerm-Large congelado no
separan linealmente los diagnósticos castellanos finos al nivel
necesario para superar L3 $\sim 0{,}15$. Ningún ajuste del lado
textual puede extraer información discriminativa que las
\emph{features} visuales no contienen. En segundo lugar, la
adaptación LoRA visual de v3 (últimos cuatro bloques,
$1{,}05$\,M parámetros) eleva el techo LP sólo $+0{,}01$ en
cada nivel, del orden de una imagen de test, dentro del ruido.
La LoRA confinada a los últimos cuatro bloques constituye
capacidad insuficiente para reesculpir las \emph{features} de
PanDerm-Large en el régimen \emph{fine-grained} castellano del
archivo. El veredicto técnico es \emph{``image LoRA did not
help''} (no destruyó el preentrenamiento, el LP nunca cayó, pero
tampoco lo mejoró sustancialmente).

El mecanismo por el cual el \emph{zero-shot} de v3 cae respecto
a v2 sin que el LP suba es interpretable: el entrenamiento
conjunto desplaza los embeddings visuales hacia el espacio
contrastivo de captions ricas y los aleja del espacio de
plantillas de nombre de clase que el \emph{zero-shot} usa
internamente. La evidencia directa es que L1 \emph{zero-shot}
\emph{accuracy}@1 cae $0{,}426$ (e$2$) $\to 0{,}287$ (e$7$)
mientras el \emph{train loss} sigue descendiendo a $0{,}29$, lo
que indica sobreajuste sobre el régimen \emph{rich-caption}.

El diagnóstico se verifica con las dos rutas candidatas de
solución, ya reportadas en la tabla~\ref{tab:spanderm-zs}: la
capacidad arquitectónica adicional (v4) confirma ser la palanca
efectiva, mientras que la escala externa genérica (v5) no replica
la mejora. En corpora clínicos pequeños de alta especificidad
lingüística, la palanca correcta para superar el techo $\sim 0{,}13$
del encoder congelado es ampliar la capacidad arquitectónica del
lado visual sobre el propio dataset, no incorporar imágenes
externas con \emph{captions} sintéticas sobre dominios afines pero
estilísticamente discordantes.

\subsection{Calibración por \emph{temperature scaling}}

Los logits \emph{zero-shot} de la rama contrastiva de Dermapixel R0
están severamente
sobreconfiados, propiedad esperada en arquitecturas CLIP-style
entrenadas con InfoNCE. La calibración con \emph{temperature
scaling} de un único parámetro escalar~\cite{guo2017} sobre los
$107$ pares de validación reduce el \emph{Expected Calibration
Error} (ECE) medido sobre los $101$ pares de test. La
tabla~\ref{tab:spanderm-calibration} sintetiza el resultado sobre
el \emph{checkpoint} v4 y la
figura~\ref{fig:spanderm-calibration} ilustra el diagrama de
fiabilidad para L3.

\begin{table}[H]
\centering
\footnotesize
\setlength{\tabcolsep}{4pt}
\renewcommand{\arraystretch}{0.95}
\caption{Calibración post-hoc por \emph{temperature scaling} sobre
la iteración contrastiva v4 de Dermapixel R0. Ajuste del escalar
$\tau$ por optimización LBFGS
sobre $107$ pares de validación; ECE medido sobre $101$ pares de
test. La calibración sobre v5 no converge (la optimización LBFGS
diverge sobre el espacio embebido \emph{image-LoRA} expandido y
$\tau \to 0$), por lo que queda registrada como limitación
metodológica abierta.}
\label{tab:spanderm-calibration}
\begin{tabular}{lrrr}
\hline
Nivel & ECE antes  & ECE después & $\tau$       \\
\hline
L3 (v4) & $0{,}7262$ & $0{,}0213$  & $0{,}4165$ \\
\hline
\end{tabular}
\end{table}

\begin{figure}[H]
\centering
\includegraphics[width=0.88\textwidth]{fig_spanderm_calibration.png}
\caption{Diagrama de fiabilidad sobre L3 antes y después de
\emph{temperature scaling} ($\tau = 0{,}4165$) sobre la iteración
contrastiva v4 de Dermapixel R0. ECE desciende de $0{,}7262$ a
$0{,}0213$ ($-97\%$). La
calibración escalar de un parámetro restaura la interpretabilidad
probabilística de los logits sin alterar el \emph{ranking} entre
clases.}
\label{fig:spanderm-calibration}
\end{figure}

La reducción es del $-97\%$ sobre L3 con un único parámetro
escalar~\cite{guo2017}. La implicación converge con la sensibilidad
de DermLIP documentada en sección~\ref{sec:zs-results}: sin
calibración post-hoc, las probabilidades de un sistema CLIP-style
\emph{zero-shot} castellano no son interpretables como confianza
clínica, por lo que cualquier despliegue debe incorporar este
escalado como capa final, ajustado sobre el archivo de destino. La
extensión del procedimiento a v5 no converge ($\tau \to 0$ sobre el
espacio \emph{image-LoRA} expandido); la limitación no afecta al
resultado principal asociado a v4 y queda registrada como
observación metodológica para iteraciones futuras.

\subsection{Trade-off entre escala externa y especificidad de
dominio: ablación v4 vs.~v5}
\label{sec:spanderm-escala-especificidad}

La asimetría del efecto v4$\to$v5 se ilumina al descomponer por
régimen de inferencia. Bajo régimen visual puro \emph{H1} ---logits
sobre la similitud entre embedding visual de test y prototipos
visuales castellanos--- la transición mejora la \emph{accuracy}@1
L3 de $0{,}168$ a $0{,}238$ ($+41{,}2\%$ relativo) y el techo
\emph{Linear Probing} L2 supervisado sube $+54{,}4\%$ relativo en
BAcc, lo que confirma que el dataset externo fortalece el encoder
visual aisladamente. Pero bajo régimen híbrido \emph{H6}
---combinación con el lado textual sobre plantillas ricas--- la
cifra global cae.

La interpretación que reconcilia ambas observaciones es que las
\emph{captions} sintéticas generadas sobre el dataset externo
genérico (PAD-UFES-20~\cite{padufes}, DermNet~\cite{dermnet}) no
preservan el registro discursivo del archivo
Dermapixel~1.0: el modelo aprende a alinear los \emph{embeddings}
de imagen con un espacio textual mixto que diluye la
\emph{fingerprint} clínica castellana sobre la que se calibra el
régimen \emph{H6}. La pérdida de alineación supera, en términos
de \emph{accuracy}@1 L3, la ganancia visual aislada. El hallazgo
es metodológicamente convergente con el documentado en
sección~\ref{sec:zs-results} sobre Derm1M: el enriquecimiento de
\emph{prompts} con material léxico no alineado con el dataset de
preentrenamiento empeora la AUROC \emph{zero-shot}, en lugar de
mejorarla por incremento de información léxica. La regularidad
empírica que ambos resultados sugieren es que la alineación
cross-modal específica de un dominio clínico de alta
especificidad lingüística constituye un recurso frágil: pequeñas
perturbaciones estilísticas del lado textual pueden trasladar la
distribución conjunta lo suficiente como para erosionar la
discriminación final aunque las representaciones marginales
mejoren.

Una segunda lectura cuantitativa refuerza el diagnóstico. La
ampliación del dataset en v5 ($854 \to 15\,982$ imágenes, factor
$\approx 19\times$) no rellena la cola larga estructural del
nivel L3: las $327$ clases con $\leq 5$ imágenes en el
subconjunto castellano original mantienen su cardinalidad tras
la incorporación de las imágenes externas, dado que el
vocabulario externo no se mapea sobre las entradas L3 finas del
archivo Taberner. Las ocho clases L3 críticas $\emph{deepened}$
en v5 ---seis tumorales más dermatitis atópica y urticaria---
ganan cobertura significativa de muestras, pero el grueso de la
cola larga permanece inactiva. La implicación operativa para el
diseño de futuras iteraciones es directa y se desarrolla en
sección~\ref{sec:linea-active-learning}: superar
\emph{zero-shot} L3 $\sim 0{,}25$ sobre vocabulario castellano
fino requiere anotación experta dirigida \emph{in-domain} sobre
las clases con peor cobertura, no escala genérica
indiscriminada.

Una ablación independiente con un modelo fundacional inglés grande
preentrenado (DermLIP v2 sobre $\approx 403\,000$ pares
Derm1M~\cite{derm1m}) confirma cuantitativamente este
diagnóstico: el \emph{wrapper} inglés con traducción de etiquetas
alcanza \emph{accuracy}@1 L3 $= 0{,}2871$ en régimen \emph{H6}
comparable, frente a $0{,}2475$ de la iteración v4 de Dermapixel R0
($+3{,}96$\,pp), y $0{,}3861$ en su mejor régimen de plantillas
más texto rico ($+13{,}86$\,pp), evidencia consistente con que la
palanca limitante era la escala de la alineación cross-modal y no
la formulación textual ni el codificador visual. El análisis
completo, junto con la cuantificación del eje idioma puro sobre el
propio \emph{checkpoint} v4 ($0{,}2475$ castellano frente a
$0{,}2277$ inglés bajo \emph{H6}) y el \emph{trade-off} operativo
entre ambas estrategias (rendimiento bruto frente a privacidad,
especificidad lingüística y autonomía de despliegue), se
desarrolla en el Anexo~\ref{cap:anexoi}.

\subsection{Diagnóstico cuantitativo del techo experimental
sobre L2}

La tabla~\ref{tab:spanderm-l2-ceiling} sitúa la iteración
contrastiva v2 de Dermapixel R0 en
\emph{zero-shot} L2 frente a las dos referencias internas
documentadas previamente en el capítulo sobre el mismo conjunto
de test ($N_{\mathrm{test}} = 36$ en
sección~\ref{sec:dermapixel-results} y
sección~\ref{sec:dermapixel-spanderm-v0}).

\begin{table}[H]
\centering
\footnotesize
\setlength{\tabcolsep}{4pt}
\renewcommand{\arraystretch}{0.95}
\caption{Rama contrastiva de Dermapixel R0 \emph{zero-shot} L2
frente a las dos
referencias internas del capítulo sobre el mismo nivel ontológico.
Las tres referencias evalúan sobre el conjunto de test de
DermapixelAI~1.0; la cifra de la iteración v2 corresponde al
test de $101$ imágenes del \emph{split} per-case del experimento
contrastivo y las dos restantes al test fijo
$N_{\mathrm{test}} = 36$ de la sección~\ref{sec:dermapixel-results}.
La comparación es indicativa por la diferencia muestral.}
\label{tab:spanderm-l2-ceiling}
\adjustbox{max width=\textwidth}{%
\begin{tabular}{llrr}
\hline
Sistema                                                 & Régimen               & Acc@1 L2          & Observaciones \\
\hline
LP PanDerm Large (referencia, tabla~\ref{tab:dermapixel-lp}) & Supervisado L2 cerrado & $0{,}250$        & $38$ clases, $N_{\mathrm{test}} = 36$ \\
Dermapixel R0 cabeza L2 (tabla~\ref{tab:dermapixel-spanderm-l2}) & Supervisado L2 cerrado & $0{,}250$        & $38$ clases, $N_{\mathrm{test}} = 36$ \\
\textbf{Dermapixel R0 contrastivo v2}                    & \emph{Zero-shot} open  & $\mathbf{0{,}133}$ & $23$ clases efectivas, $N_{\mathrm{test}} = 101$ \\
\hline
\end{tabular}}
\end{table}

La lectura cuantitativa es coherente con la jerarquía esperada
entre los tres regímenes: la cabeza supervisada cerrada con LoRA
sobre el mismo encoder (Dermapixel R0, cabeza L2) alcanza Acc@1
$0{,}250$
sobre L2, mientras que la rama contrastiva \emph{zero-shot open-class}
(iteración v2) alcanza $0{,}133$ sobre un vocabulario reducido
a las $23$ clases efectivas en el test del \emph{split}
per-case del experimento. La diferencia operativa entre ambos
regímenes es la capacidad de la rama contrastiva de aceptar consultas
textuales libres en castellano sin necesidad de fijar el
vocabulario por anticipado, propiedad explotada en el módulo
RAG castellano del prototipo (Anexo~\ref{cap:anexoh}).

\subsection{Síntesis: progresión v1$\to$v4 y ablación negativa v5}

La tabla~\ref{tab:spanderm_iteraciones} sintetiza la trayectoria
completa de la rama contrastiva de Dermapixel R0 sobre el conjunto
de test
de DermapixelAI~1.0, articulando para cada iteración la
configuración LoRA, el régimen de inferencia, las dos cifras de
referencia ($L3$ \emph{accuracy}@1 \emph{zero-shot} y lift de
recuperación cross-modal $i2t$ R@$5$) y la lectura cualitativa
asociada.

\begin{table}[H]
\centering
\footnotesize
\setlength{\tabcolsep}{3pt}
\renewcommand{\arraystretch}{1.05}
\caption{Comparativa de las cinco iteraciones del entrenamiento
contrastivo de Dermapixel R0 sobre el dataset DermapixelAI~1.0. Régimen evaluativo
\emph{zero-shot} L3 sobre $101$ imágenes de test ($54$ clases L3
representadas), \emph{H6 hybrid} v$30$\_r$70$ (combinación
$30\%$ prototipos visuales $+$ $70\%$ texto rico) aplicado de
forma homogénea a las cinco iteraciones. v4 queda a $0{,}0025$ del
objetivo fijado a priori de $0{,}25$ acc@1 L3 (dentro del ruido de
muestreo) mediante expansión arquitectónica
(\emph{image-LoRA} $r=32$ sobre $12$ bloques $+$ \emph{ensemble}
híbrido texto-visual). v5 confirma negativamente que la escala
externa genérica con \emph{captions} sintéticas no supera v4. Para
v1--v3 los prototipos visuales derivan del espacio embebido del
\emph{checkpoint} sin LoRA visual; v4--v5 emplean prototipos
sobre el espacio embebido con \emph{image-LoRA} expandido. La
comparativa permanece interpretable porque los prototipos derivan
del propio espacio embebido del \emph{checkpoint} evaluado.
Provenance: commit \texttt{d0190c3} del repositorio
\texttt{spanderm-clip}. Las cifras Linear Probing BAcc L2 absolutas
para v4 ($0{,}1860$) y v5 ($0{,}2871$) proceden de los ficheros
JSON \texttt{eval\_v5/linear\_probe\_v4ref.json} y
\texttt{eval\_v5/linear\_probe\_v5.json}; el delta relativo
$+54{,}4\%$ queda confirmado.}
\label{tab:spanderm_iteraciones}
\adjustbox{max width=\textwidth}{%
\begin{tabular}{lllrrll}
\hline
Versión & LoRA texto & LoRA visual & acc@1 L3 & Lift i2t R@$5$ & LP BAcc L2 & Notas \\
\hline
v1 & ---                       & ---                          & $0{,}1782$          & ---                       & ---                                 & Baseline e5 \emph{full FT} \\
v2 & $r=16$, $\alpha=32$       & ---                          & $0{,}1980$          & $\mathbf{14{,}0\times}$   & ---                                 & Pico retrieval rich \\
v3 & $r=16$, $\alpha=32$       & $r=16$, últ.~$4/24$           & $0{,}1980$          & ---                       & ---                                 & Image-LoRA infradotado \\
v4 & $r=32$, $\alpha=64$       & $r=32$, últ.~$12/24$          & $\mathbf{0{,}2475}$ & $12{,}4\times$            & $0{,}1860$   & A $0{,}0025$ del objetivo \\
v5 & $r=32$, $\alpha=64$       & $r=32$, últ.~$12/24$          & $0{,}2079$          & $10{,}8\times$            & $0{,}2871$   & FLAT: escala externa ($+54{,}4\%$ relativo sobre v4) \\
\hline
\end{tabular}}

\smallskip
\footnotesize
\textit{Nota histórica.} Los valores \emph{rich-only baseline} (un
\emph{caption} clínico por clase $\sim 30$--$80$ palabras
castellano) fueron $0{,}1287$ (v1), $0{,}1386$ (v2) y $0{,}1089$
(v3) en el régimen \texttt{baseline\_5\_templates} original de la
evaluación previa. La transición al régimen \emph{H6 hybrid}
homogéneo presentado en esta tabla permite la comparativa directa
cross-\emph{checkpoint} mediante el mismo procedimiento
evaluativo (prototipos visuales $30\%$ $+$ texto rico $70\%$). El
detalle metodológico se desarrolla en
Anexo~\ref{cap:anexog}.
\normalsize
\end{table}

La tabla~\ref{tab:spanderm_iteraciones} articula los tres ejes de
mejora del experimento. El eje de capacidad arquitectónica queda
validado por la transición v3$\to$v4, que confirma la hipótesis
formulada tras la medición del LP en v3: el cuello de botella
reside en la capacidad del \emph{image encoder}, no en la
formulación textual. El eje de escala externa queda
\emph{flat-to-negative} en v5, mecanísticamente discutido en la
sección~\ref{sec:spanderm-escala-especificidad} y convergente con
la sensibilidad de DermLIP al tokenizador
(sección~\ref{sec:zs-results}). El eje de supervisión conceptual
---anotaciones del \emph{Seven-Point Checklist}~\cite{kawahara2019}
o del Grupo~A (Anexo~\ref{cap:anexof})--- queda abierto como línea
prioritaria de aprendizaje activo dirigido
(sección~\ref{sec:linea-active-learning}).

En síntesis, el experimento valida la utilidad de PanDerm Large
para tareas castellanas más allá del LP/FT supervisado (v4 a
$0{,}0025$ del umbral objetivo L3 $0{,}25$ sobre vocabulario
castellano fino), disocia experimentalmente capacidad
arquitectónica de escala de datos ---palanca efectiva la primera
(v3$\to$v4), \emph{flat-to-negative} la segunda (v4$\to$v5)--- y
establece una metodología reproducible para evaluaciones
CLIP-style sobre archivos privados de pequeña escala (detallada en
el Anexo~\ref{cap:anexoi}). El sistema validado en v4 es candidato
a integrarse como módulo de búsqueda textual castellano del
prototipo DermApIxel (Anexo~\ref{cap:anexoh}), con la calibración
escalar $\tau = 0{,}4165$ como capa final.

\section{Interpretabilidad del SAE: alineamiento con conceptos
dermatoscópicos clásicos}
\label{sec:sae-derm7pt}

¿Las representaciones internas del modelo se corresponden con los
conceptos que usa un dermatólogo? La pregunta es la base de la
interpretabilidad clínica. Se halla que las \emph{features} del
autoencoder disperso se alinean con los criterios dermatoscópicos
clásicos, aunque forzar un cuello de botella conceptual cuesta
exactitud; los experimentos sobre Derm7pt lo muestran.

La línea de interpretabilidad mediante Sparse Autoencoders (SAE)
descrita en el Anexo~\ref{cap:anexof} introduce un SAE Large de
$16\,384$ \emph{features} entrenado sobre las activaciones de
PanDerm Large y valida la base conceptual del modelo sobre los
$48$ conceptos visuales de SkinCon~\cite{skincon}. La presente
sección extiende esa validación al dataset
Derm7pt~\cite{kawahara2019} ($1\,011$ lesiones dermatoscópicas) y
explota por primera vez en este trabajo las anotaciones del
\emph{Seven-Point Checklist} (siete criterios diagnósticos
estructurados por imagen). El diseño contempla tres experimentos
ejecutados sobre DGX Spark GB10 el 2026-06-02: E1 alineamiento
\emph{single-feature} entre las $16\,384$ activaciones del SAE y
los siete criterios binarizados, E2 \emph{Concept Bottleneck Model}
(CBM) jerárquico con cuello de botella sobre los siete conceptos
para predicción melanoma/no-melanoma, y E3 cruce con los $16$
conceptos dermatoscópicos del Grupo A propuestos a la
Dra.~R.~Taberner para anotación sobre DermapixelAI~1.0. Los tres
experimentos comparten el dataset completo Derm7pt
($N = 1\,011$, $24{,}5\%$ melanomas), con \emph{splits} oficiales
\emph{train}~$=413$, \emph{val}~$=203$, \emph{test}~$=395$
($100$ melanomas en test) cuando aplica división. La activación
del SAE sobre el dataset muestra $3\,214$ \emph{features} activas en
más del $5\%$ de las imágenes, $2\,978$ \emph{features}
inactivas (\emph{dead}) y $2\,717$ \emph{features} muy activas
($>50\%$), distribución coherente con la sparsidad nominal del
SAE.

El experimento E1 evalúa, para cada uno de los siete criterios
binarizados como \emph{present}/\emph{absent}, la AUROC efectiva
$\max(\mathrm{AUROC}, 1-\mathrm{AUROC})$ de cada una de las
$16\,384$ \emph{features} individuales como predictor binario
del criterio. La tabla~\ref{tab:sae-d7pt-e1} reporta el AUROC
efectivo de la mejor \emph{feature} por criterio
(\emph{top-1}) y el promedio sobre las diez mejores
(\emph{top-10}). Todas las cifras se calculan sobre el dataset
completo Derm7pt ($N = 1\,011$).

\begin{table}[H]
\centering
\footnotesize
\setlength{\tabcolsep}{4pt}
\renewcommand{\arraystretch}{0.95}
\caption{Alineamiento \emph{single-feature} entre las
$16\,384$ \emph{features} del SAE Large y los siete criterios
del \emph{Seven-Point Checklist} sobre Derm7pt
($N = 1\,011$). AUROC efectiva
$\max(\mathrm{AUROC}, 1-\mathrm{AUROC})$ de la mejor
\emph{feature} individual (\emph{top-1}) y promedio sobre las
diez mejores (\emph{top-10}). Mejor cifra agregada en negrita.}
\label{tab:sae-d7pt-e1}
\begin{tabular}{lrrrr}
\hline
Criterio                 & $N_{+}$ & $N_{-}$ & AUROC \emph{top-1} & AUROC \emph{top-10} \\
\hline
pigment\_network         & $611$ & $400$ & $0{,}754$ & $0{,}737$ \\
streaks                  & $358$ & $653$ & $0{,}688$ & $0{,}684$ \\
pigmentation             & $423$ & $588$ & $0{,}700$ & $0{,}657$ \\
regression\_structures   & $253$ & $758$ & $0{,}687$ & $0{,}673$ \\
dots\_and\_globules      & $782$ & $229$ & $0{,}681$ & $0{,}663$ \\
blue\_whitish\_veil      & $195$ & $816$ & $0{,}746$ & $0{,}721$ \\
vascular\_structures     & $188$ & $823$ & $\mathbf{0{,}823}$ & $\mathbf{0{,}782}$ \\
\hline
\end{tabular}
\end{table}

Las siete \emph{features} más alineadas individualmente alcanzan
AUROC efectivo en el intervalo $[0{,}681,\,0{,}823]$, con el
máximo sobre \emph{vascular structures} ($0{,}823$) y mínimo sobre
\emph{dots and globules} ($0{,}681$). La estabilidad entre el
\emph{top-1} y el \emph{top-10} (diferencia $\leq 0{,}043$ en
valor absoluto) indica que el alineamiento no se concentra en una
única \emph{feature} aislada sino que el SAE codifica los
conceptos mediante un pequeño subconjunto coherente de
activaciones.

El experimento E2 construye un CBM jerárquico de dos etapas:
una primera cabeza de regresión logística sobre las
$16\,384$ activaciones del SAE estima la probabilidad de cada
uno de los siete criterios; un meta-clasificador logístico toma
estas siete estimaciones como vector de entrada y predice
melanoma/no-melanoma. La comparativa con el sondeo lineal directo
sobre las $16\,384$ \emph{features} sin cuello de botella se
reporta en la tabla~\ref{tab:sae-d7pt-cbm}.

\begin{table}[H]
\centering
\footnotesize
\setlength{\tabcolsep}{4pt}
\renewcommand{\arraystretch}{0.95}
\caption{Comparativa entre sondeo lineal directo sobre las
$16\,384$ \emph{features} del SAE Large y CBM jerárquico con
los siete criterios \emph{Seven-Point Checklist} como cuello de
botella. Conjunto de test Derm7pt
($N_{\mathrm{test}} = 395$, $100$ melanomas). AUROC en
negrita.}
\label{tab:sae-d7pt-cbm}
\begin{tabular}{lrrrrr}
\hline
Arquitectura                                & Acc@1     & BAcc      & AUROC              & W-F1      & Kappa     \\
\hline
Sondeo lineal directo (SAE $\to$ melanoma)  & $0{,}841$ & $0{,}735$ & $\mathbf{0{,}890}$ & $0{,}829$ & $0{,}527$ \\
CBM (SAE $\to$ $7$ conceptos $\to$ melanoma) & $0{,}825$ & $0{,}681$ & $\mathbf{0{,}832}$ & $0{,}801$ & $0{,}440$ \\
\hline
\end{tabular}
\end{table}

El paso de las activaciones crudas del SAE a través del cuello de
botella conceptual reduce la AUROC binaria melanoma de $0{,}890$ a
$0{,}832$ ($-5{,}8$\,pp), pérdida que documenta el compromiso entre
interpretabilidad y precisión. Los pesos del meta-clasificador
identifican \emph{blue whitish veil} y \emph{regression structures}
como los conceptos dominantes para la decisión melanoma, jerarquía
coherente con el peso clínico que el \emph{Seven-Point
Checklist}~\cite{kawahara2019} asigna a ambos criterios (criterios
mayores de $2$ puntos). El AUROC binario por concepto de la cabeza
intermedia y los pesos completos del meta-clasificador figuran en el
Anexo~\ref{cap:anexof} (tabla~\ref{tab:sae-d7pt-cbm-coefs}).

El experimento E3 cruza los $16$ conceptos dermatoscópicos del
Grupo~A propuestos a la Dra.~R.~Taberner para anotación sobre las
$49$ imágenes dermatoscópicas de DermapixelAI~1.0
(sección~\ref{sec:dermapixel}) con los siete criterios del
\emph{Seven-Point Checklist} sobre Derm7pt. Para cada concepto
Rosa con \emph{mapping} a un criterio Derm7pt (directo o por
subcategoría) se reporta la AUROC efectiva de la mejor
\emph{feature} SAE sobre la binarización correspondiente. El
resultado completo se reporta en el Anexo~\ref{cap:anexof}
(tabla~\ref{tab:sae-d7pt-cross-rosa}). De los $16$ conceptos del
Grupo~A, $11$ admiten \emph{mapping} operativo a un criterio del
\emph{Seven-Point Checklist}, y de ellos $7$ quedan cubiertos por al
menos una \emph{feature} SAE con AUROC efectivo $\geq 0{,}70$ ---con
un máximo de $0{,}926$ sobre vasos en horquilla, pese a sólo $15$
positivos sobre $1\,011$ imágenes, lo que ilustra que la
representación SAE captura conceptos poco frecuentes sin
contaminación por la distribución mayoritaria. Los $5$ conceptos sin
\emph{mapping} (patrón empedrado, paralelos surcos/crestas, lagunas
vasculares, vasos en corona) corresponden a estructuras no
contempladas en la escala original y definen el contenido específico
de la anotación clínica colaborativa con la Dra.~R.~Taberner.

\subsection*{Ajuste fino multitarea: cabezas conceptuales y
melanoma}
\label{sec:sae-d7pt-e4}

El experimento E4 replica conceptualmente el protocolo multitarea
descrito por Kawahara et~al.~\cite{kawahara2019} sustituyendo la
red convolucional original por el codificador fundacional PanDerm
Large. El \emph{fine-tuning} se realiza mediante adaptación de
bajo rango (LoRA, $r = 16$) aplicada exclusivamente a los dos
últimos bloques del codificador (\emph{blocks.22-23}), de forma
idéntica al protocolo de Dermapixel R0 reportado en la
sección~\ref{sec:dermapixel-spanderm-v0}. Sobre la
representación adaptada se conectan ocho cabezas paralelas: una
cabeza multi-clase por cada uno de los siete criterios del
\emph{Seven-Point Checklist} y una cabeza binaria para
melanoma/no-melanoma. La función objetivo es la suma
equiponderada de entropías cruzadas sobre las ocho cabezas,
$\mathcal{L} = \sum_{i=1}^{7} \mathrm{CE}(\hat{c}_i, c_i) +
\mathrm{CE}(\hat{m}, m)$, con \emph{class weight} balanceado por
clase. La configuración de optimización emplea AdamW con
\emph{learning rate} diferenciado
($\eta_{\mathrm{heads}} = 10^{-3}$,
$\eta_{\mathrm{LoRA}} = 5 \times 10^{-4}$) durante $15$ épocas,
con selección del \emph{checkpoint} de mejor AUROC melanoma sobre
validación (época $13$, AUROC val $= 0{,}861$). El número total
de parámetros entrenables es $0{,}56$\,M ($0{,}18\%$ del total
del codificador), magnitud equivalente a la de Dermapixel R0.

La tabla~\ref{tab:sae-d7pt-e4-comp} compara el rendimiento de
E4 frente al sondeo lineal directo (E2 Direct LP sobre las
$16\,384$ activaciones del SAE) y al CBM jerárquico de siete
conceptos (E2 CBM) sobre el conjunto de test Derm7pt
($N_{\mathrm{test}} = 395$, $n_{\mathrm{mel}} \approx 98$). Los
intervalos de confianza al $95\%$ se estiman por \emph{bootstrap}
no paramétrico sobre el conjunto de test.

\begin{table}[H]
\centering
\footnotesize
\setlength{\tabcolsep}{4pt}
\renewcommand{\arraystretch}{0.95}
\caption{Comparativa sobre la tarea binaria melanoma/no-melanoma
entre sondeo lineal directo sobre las $16\,384$ \emph{features}
del SAE Large (E2 Direct LP), CBM jerárquico de siete conceptos
(E2 CBM) y \emph{fine-tuning} multitarea LoRA con ocho cabezas
paralelas (E4). Conjunto de test Derm7pt
($N_{\mathrm{test}} = 395$, $n_{\mathrm{mel}} \approx 98$), mismo
\emph{split} oficial. AUROC en negrita.}
\label{tab:sae-d7pt-e4-comp}
\begin{tabular}{lrrrrrr}
\hline
Método                                       & Params entr. & Acc@1     & BAcc      & AUROC              & W-F1      & Kappa     \\
\hline
E2 Direct LP (SAE $\to$ melanoma)             & $16$\,K       & $0{,}841$ & $0{,}735$ & $0{,}890$          & $0{,}829$ & $0{,}527$ \\
E2 CBM (SAE $\to$ $7$ conceptos $\to$ mel)    & $\sim 120$    & $0{,}825$ & $0{,}681$ & $0{,}832$          & $0{,}801$ & $0{,}440$ \\
\textbf{E4 FT multitarea LoRA}                & $0{,}56$\,M   & $\mathbf{0{,}841}$ & $\mathbf{0{,}804}$ & $\mathbf{0{,}891}$ & $\mathbf{0{,}843}$ & $\mathbf{0{,}590}$ \\
\hline
\end{tabular}
\end{table}

El desglose por concepto que compara el sondeo lineal separado
(E2) con las siete cabezas del \emph{fine-tuning} multitarea (E4)
figura en el Anexo~\ref{cap:anexof}
(tabla~\ref{tab:sae-d7pt-e4-concepts}): las cabezas conceptuales del
régimen multitarea pierden entre $1$ y $3$\,pp de AUROC binaria
frente al sondeo lineal separado en seis de los siete conceptos
(máximo $-2{,}7$\,pp en \emph{regression structures}), una
degradación menor que confirma que el aprendizaje conjunto de
melanoma y conceptos no compromete la calidad conceptual.

Tres observaciones cuantitativas sintetizan E4. \textit{Primero},
la cabeza melanoma del \emph{fine-tuning} multitarea alcanza
AUROC test $= 0{,}891$ (IC95\% $[0{,}856,\,0{,}927]$), valor que
supera marginalmente al sondeo lineal directo sobre las
$16\,384$ \emph{features} del SAE ($0{,}890$, $+0{,}1$\,pp,
incremento dentro del ruido del \emph{bootstrap}) y mejora el
CBM jerárquico en $+5{,}9$\,pp AUROC. La diferencia operativa
respecto al Direct LP es que E4 ofrece simultáneamente predicción
melanoma e inferencia sobre los siete criterios diagnósticos en
una sola pasada, lo que preserva la interpretabilidad explícita
sin coste discriminativo apreciable sobre la tarea binaria
primaria. \textit{Segundo}, las siete cabezas conceptuales pierden
entre $1$ y $3$\,pp de AUROC respecto al sondeo lineal separado
por concepto: el coste de compartir representación entre tareas
auxiliares y tarea principal se sitúa en
$-1{,}2$\,pp (\emph{pigment network}), $-1{,}6$\,pp (\emph{blue
whitish veil}), $-2{,}6$\,pp (\emph{vascular structures}) y
$-2{,}7$\,pp (\emph{regression structures}); dos conceptos
mejoran ligeramente bajo el régimen multitarea
(\emph{pigmentation} $+0{,}1$\,pp,
\emph{dots and globules} $+0{,}2$\,pp). Este patrón es
consistente con el comportamiento documentado del aprendizaje
multitarea y cuantifica el compromiso aceptado a cambio de un
único codificador adaptado para las ocho cabezas
simultáneas. \textit{Tercero}, el hecho de que la cabeza melanoma
no sufra penalización medible respecto al sondeo lineal
especializado sugiere que el aprendizaje multitarea no degrada la
tarea principal cuando los conceptos auxiliares son
taxonómicamente coherentes con el objetivo. El
\emph{Seven-Point Checklist} es, por construcción, un instrumento
diagnóstico de melanoma~\cite{kawahara2019}, lo que justifica
empíricamente la elección del dataset de criterios.

Cuatro lecturas sintetizan E1--E4 sin reproducir las cifras ya
reportadas. El SAE Large entrenado sobre Derm1M codifica conceptos
dermatoscópicos clásicos sin supervisión explícita (AUROC efectivo
hasta $0{,}926$ sobre vasos en horquilla) y generaliza a Derm7pt
sin reentrenamiento, lo que sostiene la línea de interpretabilidad
del Anexo~\ref{cap:anexof}. El CBM jerárquico cuantifica el
compromiso interpretabilidad--precisión ($-5{,}8$\,pp AUROC frente
al sondeo lineal directo) y ordena \emph{blue whitish veil} y
\emph{regression structures} como conceptos dominantes, jerarquía
coherente con el peso clínico del \emph{Seven-Point
Checklist}~\cite{kawahara2019}. El \emph{fine-tuning} multitarea
LoRA (E4) iguala al sondeo lineal directo sobre la tarea binaria
melanoma ($0{,}891$ frente a $0{,}890$) preservando inferencia
explícita concepto-a-concepto, a cambio de $1$--$3$\,pp en las
cabezas conceptuales. De los $11$ conceptos del Grupo~A con
\emph{mapping} a Derm7pt, $7$ quedan cubiertos por \emph{features}
SAE; los $5$ sin \emph{mapping} (empedrado, paralelos
surcos/crestas, lagunas vasculares, vasos en corona) definen el
contenido específico de la anotación clínica colaborativa
documentada como línea de continuidad en la
sección~\ref{sec:linea-derm7pt-sae}. El conjunto refuerza la
viabilidad del módulo M7 de interpretabilidad del prototipo
DermApIxel (Anexo~\ref{cap:anexoh}).

\section{Síntesis comparativa de métodos sobre DermapixelAI 1.0}
\label{sec:dermapixel-extra}
\label{sec:dermapixel-extra-sintesis}

Las secciones~\ref{sec:dermapixel-results}--\ref{sec:spanderm-clip}
evalúan codificadores específicos de dominio (PanDerm Base/Large,
DermLIP v2), vision-lenguaje médicos (BiomedCLIP, SigLIP-SO400M) y
las dos ramas de Dermapixel R0 sobre DermapixelAI~1.0. El
Anexo~\ref{cap:anexoj} añade dos ablaciones complementarias
---codificadores generalistas no médicos (DINOv2 y CLIP-L) bajo
sondeo lineal, y dos arquitecturas de inferencia alternativas
(\emph{retrieval} $k$-NN con índice FAISS y \emph{zero-shot}
jerárquico en cascada)--- cuyas mejores cifras se integran en la
síntesis siguiente. La tabla~\ref{tab:dermapixel-extra-sintesis}
consolida la BAcc y la AUROC por nivel ontológico para el conjunto
completo de métodos y permite identificar el mejor por nivel y
métrica sin dispersión entre las secciones de origen.

\begin{table}[H]
\centering
\footnotesize
\setlength{\tabcolsep}{4pt}
\renewcommand{\arraystretch}{0.95}
\caption{Síntesis de codificadores y arquitecturas evaluados sobre
DermapixelAI~1.0 (test fijo: $N = 36$ en L1/L2, $N = 28$ en L3).
La columna \emph{Ref.} remite a la sección de origen de las cifras
completas con sus intervalos de confianza. Mejor valor por columna
y nivel en negrita.}
\label{tab:dermapixel-extra-sintesis}
\begin{tabular}{llrrrrrr}
\hline
                                  &     & \multicolumn{2}{c}{L1} & \multicolumn{2}{c}{L2} & \multicolumn{2}{c}{L3} \\
Método                            & Ref.& BAcc & AUROC           & BAcc & AUROC           & BAcc & AUROC \\
\hline
PanDerm Base LP                   & §4.11 & $0{,}570$          & $0{,}697$
                                          & $0{,}221$          & $0{,}707$
                                          & $0{,}132$          & $0{,}735$ \\
PanDerm Large LP                  & §4.11 & $\mathbf{0{,}735}$ & $\mathbf{0{,}796}$
                                          & $0{,}265$          & $0{,}793$
                                          & $0{,}184$          & $0{,}813$ \\
DermLIP v2 LP                     & §4.11 & ---                & ---
                                          & ---                & $0{,}860$
                                          & ---                & --- \\
DINOv2 ViT-L LP                   & Anx.~\ref{cap:anexoj} & $0{,}537$          & $0{,}650$
                                          & $0{,}250$          & $0{,}718$
                                          & $0{,}211$          & $0{,}794$ \\
CLIP-L-336 LP                     & Anx.~\ref{cap:anexoj} & $0{,}556$          & $0{,}658$
                                          & $0{,}279$          & $\mathbf{0{,}826}$
                                          & $0{,}184$          & $\mathbf{0{,}827}$ \\
Dermapixel R0 LoRA                & §4.14 & ---                & ---
                                          & $\mathbf{0{,}363}$ & $0{,}855$
                                          & ---                & --- \\
FT parcial denso L1$+$L2$+$L3      & §4.11 & $0{,}622$          & $0{,}715$
                                          & $0{,}324$          & $0{,}852$
                                          & $0{,}184$          & $0{,}804$ \\
$k$-NN FAISS PanDerm Large         & Anx.~\ref{cap:anexoj} & $0{,}686$          & $0{,}748$
                                          & $0{,}265$          & $0{,}700$
                                          & $\mathbf{0{,}210}$ & $0{,}719$ \\
\hline
\end{tabular}
\end{table}

El mejor método sobre DermapixelAI~1.0 depende del nivel
ontológico y de la métrica. Sobre L1, el sondeo lineal sobre
PanDerm Large encabeza simultáneamente BAcc ($0{,}735$) y
AUROC ($0{,}796$); sobre L2, Dermapixel R0 LoRA encabeza la BAcc
($0{,}363$) y CLIP-L-336 LP encabeza la AUROC ($0{,}826$); sobre
L3, el \emph{retrieval} $k$-NN FAISS sobre embeddings de PanDerm
Large encabeza la BAcc ($0{,}210$) y CLIP-L-336 LP encabeza la
AUROC ($0{,}827$). La discriminación binaria melanoma/no-melanoma
sobre Derm7pt (sección~\ref{sec:sae-derm7pt}, fuera de
DermapixelAI~1.0) alcanza AUROC $0{,}891$ con el \emph{fine-tuning}
multitarea LoRA E4 sobre las activaciones del SAE. La ausencia de
un método único óptimo refuerza la lectura interpretativa de la
sección~\ref{sec:disc-especializacion}: la elección del
codificador y del clasificador debe condicionarse al nivel
ontológico objetivo y al criterio operativo predominante
(discriminación robusta, ranking probabilístico o recuperación
para soporte explicativo).


\section{Aplicación de seguridad: ensemble para detección de melanoma}
\label{sec:ensemble-results}

Más allá de la comparación de modelos individuales, resulta
relevante explorar si la complementariedad entre arquitecturas puede
explotarse para maximizar la sensibilidad diagnóstica. Esta cuestión
es especialmente importante en melanoma, donde el coste clínico de
un falso negativo es potencialmente muy superior al de un falso
positivo. Por ello se evalúa un \emph{ensemble} orientado
explícitamente a maximizar el \emph{recall} de melanoma sobre
HAM10000.

Se combinan tres clasificadores independientes sobre el
\emph{split} de test de HAM10000 ($N = 1\,232$ imágenes,
$70$ melanomas):
M1 (PanDerm Large \emph{fine-tuned} sobre HAM10000), SigLIP-Large
SO400M con sondeo lineal sobre los \emph{embeddings} de
$1\,152$ dimensiones, y M7 (clasificador unificado sobre el
\emph{merged43} del dataset armonizado, descrito en el
Anexo~\ref{cap:anexoh}).

\begin{table}[H]
\centering
\footnotesize
\setlength{\tabcolsep}{4pt}
\renewcommand{\arraystretch}{0.95}
\caption{Ensemble sobre HAM10000 ($N = 1\,232$, $70$ melanomas).
\emph{Recall} \emph{top-3}: la clase melanoma aparece entre las
tres primeras del ranking. FP\,mel: falsos positivos melanoma
(imágenes no-melanoma marcadas).}
\label{tab:ensemble}
\begin{tabular}{lrrrrrr}
\hline
Configuración                   & Acc           & BAcc          & AUROC         & Mel rec.\ top-1 & Mel rec.\ top-3 & FP\,mel \\
\hline
M1 (PanDerm FT)                 & $0{,}919$     & $0{,}813$     & $0{,}986$     & $0{,}714$ (50/70) & $0{,}971$ (68/70) & $36$ \\
SigLIP-L LP                     & $0{,}900$     & $0{,}776$     & $0{,}971$     & $0{,}614$ (43/70) & $0{,}986$ (69/70) & $31$ \\
Avg M1 + SigLIP                 & $0{,}917$     & $0{,}836$     & $0{,}986$     & $0{,}686$         & $0{,}986$         & $32$ \\
MaxMel M1 + SigLIP              & $0{,}909$     & $0{,}832$     & $0{,}985$     & $\mathbf{0{,}771}$ & $\mathbf{1{,}000}$ & $50$ \\
OR top-3 M1 + SigLIP            & ---           & ---           & ---           & ---               & $\mathbf{1{,}000}$ (70/70) & $885$ \\
\hline
\end{tabular}
\end{table}

La combinación OR \emph{top-3} de M1 y SigLIP alcanza un
\emph{recall} de melanoma del $100\%$ ($70$ de $70$), pero a costa
de $885$ falsos positivos sobre las $1\,232$ imágenes del test: la
sensibilidad máxima se obtiene marcando como sospechosa cerca del
$72\%$ del conjunto, lo que reduce drásticamente la precisión y
acota su utilidad a un escenario de \emph{cribado}, nunca de
decisión diagnóstica.

El principal hallazgo de esta sección, sin embargo, no es ese
$100\%$ ---una cifra trivialmente alcanzable ampliando el margen de
sospecha de cualquier clasificador--- sino la evidencia de
complementariedad entre una arquitectura especializada (M1, PanDerm
\emph{fine-tuned} \emph{in-domain}) y una generalista (SigLIP con
sondeo lineal): ambas recuperan melanomas distintos. Los dos
melanomas que M1 omite en \emph{top-3} (ISIC\_0024886.jpg e
ISIC\_0030360.jpg) sí los recupera SigLIP, y el único melanoma que
SigLIP no sitúa en \emph{top-3} sí lo recupera M1. Este resultado
indica que los errores de ambos modelos no están completamente
correlacionados ---condición necesaria para que una estrategia de
\emph{ensemble} aporte un beneficio clínico real, y no una mera
redundancia--- y sugiere que cada arquitectura explota
representaciones parcialmente diferentes de la lesión. M7, evaluado
sobre el subconjunto de $198$ imágenes con predicción simultánea de
los tres modelos ($12$ melanomas), no recupera melanomas adicionales
sobre el par, pero aporta cobertura sobre las $43$ clases
ontológicas del dataset armonizado.

Tres caveats acompañan la interpretación de este resultado: (i)~el
\emph{recall} del $100\%$ es \emph{top-3}, no \emph{top-1};
(ii)~el tamaño muestral de melanomas es $N = 70$, con intervalos
de confianza al $95\%$ amplios; (iii)~no existe validación
prospectiva sobre cohorte hospitalaria. El resultado se interpreta
como prueba de concepto para un escenario de cribado oportunista,
no como sustituto del juicio clínico.
